\documentclass[11pt,a4paper]{article}
\usepackage[T1]{fontenc}
\usepackage[utf8]{inputenc}
\usepackage[margin=25mm]{geometry}
\usepackage{amsmath,amssymb,booktabs,longtable,tabularx}
\usepackage{graphicx}
\usepackage{mathrsfs}
\usepackage[hidelinks]{hyperref}
\urlstyle{tt}
\hypersetup{pdftitle={SIDIS through NLO: calculation methods and numerical results},
 pdfauthor={SIDIS calculation workflow},
 pdfsubject={Six channels with reverse unitarity, Kira and SubTropica}}
\setlength{\parindent}{0pt}
\setlength{\parskip}{0.55em}
\setlength{\emergencystretch}{2em}
\setcounter{tocdepth}{2}
\allowdisplaybreaks[1]
\newcommand{\dd}{\mathrm{d}}
\newcommand{\eps}{\epsilon}
\newcommand{\MSbar}{\overline{\mathrm{MS}}}
\newcommand{\Fhat}{\widehat F}
\newcommand{\What}{\widehat W}
\newcommand{\NLO}{\mathrm{NLO}}
\newcommand{\LO}{\mathrm{LO}}
\newcommand{\PDF}{\mathrm{PDF}}
\newcommand{\FF}{\mathrm{FF}}
\newcommand{\ct}{\mathrm{CT}}
\newcommand{\cut}{\mathrm{cut}}
\newcommand{\obs}[1]{\underline{#1}}
\newcommand{\Li}{\operatorname{Li}}
\newcommand{\GeV}{\mathrm{GeV}}
\newcommand{\pb}{\mathrm{pb}}

\title{The SIDIS calculation through NLO\\
\large Six channels with reverse unitarity, Kira and SubTropica}
\author{Calculation-method report for the local SIDIS workflow}
\date{7 September 2026}

\begin{document}
\maketitle

This report explains the accepted calculation in \path{scripts/SIDIS/}.
The outputs are the symbolic partonic structure functions
$\Fhat_{1,ij}$ and $\Fhat_{2,ij}$ for all six channels.
The notation and definition of the observable follow the supplied SIDIS
paper~\cite{sidispaper}. The real integral calculation follows the
reverse-unitarity approach of the supplied standalone reference~\cite{standalone}.
Equations below state the definitions used by those sources and the programs~\cite{localcalculation},
or quote saved tool results. No new physics calculation was performed to
write this report.

The implemented routes are
\begin{align*}
\text{real:}\quad&
\text{FeynArts}\to\text{FeynCalc}\to\text{reverse unitarity}
\to\text{Kira}\to\text{SubTropica},\\
\text{virtual and self energies:}\quad&
\text{FeynArts}\to\text{FeynCalc}\to\text{Kira}\to\text{SubTropica}.
\end{align*}
These are followed by renormalization where needed, collinear factorization,
pole cancellation and the final structure-function projection.
The unchanged generated amplitudes, unintegrated tensors, Born tensors and
observable definitions are imported from the existing channel calculations.
The required real, virtual and self-energy master values are evaluated in
the new integral workflow. Package-X/PaVe values and the paper's evaluated
angular-integral tables do not supply those values.

\paragraph{How to read this report.}
The common sections define the observable, amplitude algebra, real and
virtual integration, and subtraction. Each channel then has its own
section stating its generated process, Born input where present, flavor
weights, subtraction routes and final output. The numerical sections
follow an F hat into a measured transverse-momentum bin: kinematic maps,
PDFs, fragmentation, acceptance, endpoint subtraction, integration, and
the two saved figures. The final file tables identify the full symbolic
coefficients, whose long expressions remain in executable Wolfram Language
form.

\paragraph{Scope of the calculation.}
The hard scattering is massless, unpolarized, and evaluated through NLO
at nonzero transverse momentum. Symbolic $N_c$, $N_f$ and quark charges
are retained in the production hard parts. Numerical choices such as
$N_c=3$, active flavors, a common scale and neutral-pion fragmentation
are made by the numerical consumer. This document describes the
accepted run and its stated limitations; it does not replace the
calculation by a derivation performed while writing the report.

\clearpage
\tableofcontents
\clearpage

\section{Observable, momenta and channels}
\label{sec:observable}

The process is
\begin{equation}
 \ell(l)+H(P_H)\longrightarrow \ell'(l')+h(P_h)+X,
 \qquad
 \gamma^*(q)+i(p)\longrightarrow j(k_1)+X.
\end{equation}
Here $P_H$ denotes the incoming hadron momentum, and $P_h$ the observed
hadron momentum. The first channel label is the incoming parton; the second
is the observed, fragmenting parton. Its momentum is always $k_1$.
Only the unobserved partons are integrated in the partonic hard tensor.
The calculation is unpolarized and uses the azimuth-independent $F_1,F_2$
decomposition of the paper.

The defining relations are
\begin{gather}
 D=4-2\eps,\qquad q^2=-Q^2<0,\qquad p^2=k_1^2=k_2^2=k_3^2=0,
 \label{eq:dimension}\\
 s=(p+q)^2,\qquad t=(q-k_1)^2,\qquad u=(p-k_1)^2,\\
 P_X=p+q-k_1,\qquad w=P_X^2=s_{23},\qquad
 s+t+u=-Q^2+w,\label{eq:invariants}\\
 \hat x=\frac{Q^2}{2p\cdot q}=\frac{x}{\xi},\qquad
 \hat z=\frac{p\cdot k_1}{p\cdot q}=\frac{z}{\zeta},\qquad
 p=\xi P_H,\quad k_1=P_h/\zeta.
\end{gather}
These are the paper's definitions, implemented and checked in shared S02.
The Born and virtual kinematics have $w=0$. The interior real-emission
domain recorded in the final files is
\begin{equation}
 Q^2>0,\qquad s>w>0,\qquad
 w-Q^2-s<t<-\frac{Q^2w}{s}.
 \label{eq:support}
\end{equation}
Endpoint distributions are defined separately at $w=0$.

\subsection{What is observed and what is integrated}

There are two different integrations. At the parton level, $p,q,k_1$
are external data, and the unresolved final state is integrated to
construct $\What^{\mu\nu}$. At the hadron level, the resulting hard
functions are convolved with the incoming PDF and the fragmentation
function of the observed parton. Reverse unitarity is used for the first
integration. The PDF/FF convolution and the experimental bin integration
are performed later.

The observed-particle label matters even when two rows contain the same
species. Hqq and Hqg share the Born scattering $\gamma^*q\to qg$, but
attach fragmentation to different outgoing legs. Their definition of
$t$, the tensor contractions, the fragmentation subtraction and the final
coefficient must all use their own $k_1$ assignment.

\subsection{Perturbative order and the photon convention}

At nonzero transverse momentum the leading hard parts are of order
$\alpha_s$, and this calculation keeps the correction of order
$\alpha_s^2$. This is the meaning of LO and NLO in this report.
The following table lists the generated processes; an underline marks $k_1$.
The photon is understood on the incoming side of every row.

\begin{center}
\small
\begin{tabularx}{\textwidth}{@{}l l l X@{}}
\toprule
Folder & Incoming & Born final state & Real final states at $\alpha_s^2$\\
\midrule
\texttt{Hqq} & $q_i$ & $\obs{q_i}\,g$ &
 $\obs{q_i}\,gg$; $\obs{q_i}\,q_i\bar q_i$;
 $\obs{q_i}\,q_j\bar q_j$ with $j\ne i$\\
\texttt{Hqg} & $q_i$ & $\obs{g}\,q_i$ & $\obs{g}\,q_i g$\\
\texttt{Hgq} & $g$ & $\obs{q_i}\,\bar q_i$ & $\obs{q_i}\,\bar q_i g$\\
\texttt{Hgg} & $g$ & absent & $\obs{g}\,q_f\bar q_f$\\
\texttt{Hqqbar} & $q_i$ & absent & $\obs{\bar q_i}\,q_i q_i$\\
\texttt{Hqqprime} & $q_i$ & absent & $\obs{q_j}\,q_i\bar q_j$, $j\ne i$\\
\bottomrule
\end{tabularx}
\end{center}
Hqq, Hqg and Hgq require one-loop corrections to their Born processes.
Hgg, Hqqbar and Hqqprime first contribute at the real-emission order shown;
they have no virtual contribution at this order.

The electromagnetic coupling used in the lepton cross section is kept
outside the partonic F hats. Squared quark charges remain in the hard
functions with the conventions stated for each channel. The generated
Born, real and one-loop amplitudes carry the recorded coupling degrees
$(e,g_s)$ of $(1,1)$, $(1,2)$ and $(1,3)$, respectively. The virtual
contribution is the interference with the Born amplitude. The square of
a one-loop amplitude is outside this NLO calculation.

\paragraph{Symbols in the saved files.}
The symbol \texttt{Q2} means $Q^2$, \texttt{w} means $s_{23}$, and
\texttt{SUNN} or \texttt{Nc} means $N_c$. The symbols \texttt{mu} and
\texttt{mu2} retain their source-specific scale conventions; the latter
is a squared scale. In the numerical sections,
$p_T=|\boldsymbol P_{hT}|$ is the hadron transverse momentum and $z$ is
the hadron fragmentation variable of the paper.

\section{From diagrams to two scalar contractions}
\label{sec:amplitudes}

\subsection{Generation and the reusable inputs}

FeynArts supplies the diagrams and amplitudes. FeynCalc performs the spin
and color algebra and the tensor contractions. The imported inputs retain
the generated diagram factors, incoming averages, photon-charge removal,
observed-particle tags and identical-particle bookkeeping.
The imported real interference pieces must sum exactly to the saved
unintegrated tensors. In the virtual route the Born interference must
reproduce each channel's own accepted Born contractions.

The generation programs use the massless QCD part of the FeynArts
\texttt{SMQCD} model. Requested external fields, loop order and
observed-leg ordering define each process. The core scattering sources
exclude tadpoles and external wave-function corrections from diagram
generation; external-field effects enter at renormalization. Internal
fields are restricted to the intended QCD content. FeynCalc's
\texttt{FCFAConvert} converts the amplitudes to $D$ dimensions, with
the recorded mass substitutions and coupling conventions. Original
diagrams and amplitude lists are preserved alongside the scalar inputs.

These inputs remain useful when the integration method changes. The new
import accepts the generated amplitudes, Born contractions, unintegrated
real tensors, virtual integrands and loop-independent subtraction
definitions, with their recorded identities. An evaluated old real
angular integral or an old Package-X virtual value is not used as a new
master integral.

\subsection{Interferences, spin sums and color algebra}

The square is assembled from generated diagram pairs. Diagonal pairs and
off-diagonal interferences carry the factors recorded by the pair
program. FeynCalc performs fermion spin sums, Dirac traces, polarization
sums and $SU(N_c)$ color contractions before scalar integral reduction.
Pair results can be cached separately, so the complete unsimplified
diagram sum need not be held in memory at once.

The saved color and spin measurements in the core Born records give
\begin{gather}
 C_F=\frac{(N_c-1)(N_c+1)}{2N_c},\qquad
 T_F=\frac12,\qquad C_A=N_c,\label{eq:color}\\
 n_{\mathrm{spin},q}=2,\qquad n_{\mathrm{spin},g}=D-2,\\
 A_q=\frac{1}{2N_c},\qquad
 A_g=\frac{1}{(D-2)(N_c-1)(N_c+1)}.
 \label{eq:averages}
\end{gather}
These equations quote the stored color traces and physical spin
projectors. The incoming-gluon average remains $D$ dimensional.
Its dependence on the regulator must be retained in pole-multiplied
terms.

Physical gluon polarization sums use the reference vectors prescribed
by the source, with $p$ and $k_1$ supplying the relevant lightlike
directions. Photon and gluon Ward tests are evaluated on the appropriate
generated amplitudes or contractions. A Ward test checks gauge
consistency; it does not by itself fix a spectator factor, an incoming
average or the normalization of a separate program.

\subsection{Charges and identical particles}

The source model has definite quark charges. Core normalization programs
measure the model photon vertex, remove the model charge and expose the
physical charge symbols used by the channel. In four-quark sectors,
the two photon-charge assignments are retained through their interference.
Flavor summation is performed at the documented assembly boundary.

The program distinguishes the observed tag from unobserved identical
particles. Core tagged squares obtain their weight from the generated
final-species list. The real-only convention leaves the unobserved
spectator factor outside the imported tensor. Each channel section states
which weight is already present and which enters at final assembly.

\subsection{Tensor projection and dimensional dependence}

The two contractions are
\begin{equation}
 T_g=g^{\mu\nu}\What_{\mu\nu},\qquad
 T_{pp}=p^\mu p^\nu\What_{\mu\nu}.
 \label{eq:contractions}
\end{equation}
The first uses $g^{\mu\nu}$ with the sign stated here. A sign change to a
physical-photon polarization sum must not be inserted into this definition.
Each charge and real-emission sector is retained separately until its
required weights have been applied.

The paper's partonic decomposition is
\begin{equation}
 \What^{\mu\nu}_{ij}
 =\left(-g^{\mu\nu}+\frac{q^\mu q^\nu}{q^2}\right)\Fhat_{1,ij}
 +\frac{\widetilde p^\mu\widetilde p^\nu}{p\cdot q}\Fhat_{2,ij},
 \qquad
 \widetilde p^\mu=p^\mu-q^\mu\frac{p\cdot q}{q^2}.
\end{equation}
The saved $D$-dimensional projector definitions are
\begin{align}
 \mathcal P_1^{(D)}[T_g,T_{pp}]
 &=\frac{-T_g/2+2\hat x^2 T_{pp}/Q^2}{1-\eps},\\
 \mathcal P_2^{(D)}[T_g,T_{pp}]
 &=\frac{-\hat x T_g+4(3-2\eps)\hat x^3T_{pp}/Q^2}{1-\eps}.
 \label{eq:projectors}
\end{align}
The Born terms used in subtraction retain their $D$ dependence. The final
projection uses the finite combined contractions only after their poles
have canceled.

The projector program constructs the decomposition, contracts it with
$g^{\mu\nu}$ and $p^\mu p^\nu$, and solves the linear system. It
checks the resulting projectors against the defining tensor and the
paper. Thus projectors and Born tensors share the same sign and dimension
convention. In particular, \texttt{Pg} means the $g^{\mu\nu}$
contraction; the source explicitly accounts for its sign relative to the
virtual-photon polarization-sum operation.

Keeping both scalar contractions through real, virtual and collinear
assembly lets the pole test use one normalization. The final projector
acts on finite contractions. The Born tensor inside subtraction remains
$D$ dimensional because poles multiplying its regulator dependence
contribute to the finite result.

\section{Real integration by reverse unitarity}
\label{sec:real}

\subsection{The two physical cuts}

For real emission the common loop routing is
\begin{equation}
 k_2=r,\qquad k_3=P_X-r.
\end{equation}
The two on-shell constraints become the two cut propagators of one loop.
In discontinuity notation the positive-energy cut is
\begin{equation}
 \left[\frac{1}{k^2}\right]_{\cut,+}
 \equiv\frac{\theta(k^0)}{2\pi i}
 \left(\frac{1}{k^2-i0}-\frac{1}{k^2+i0}\right)
 =\delta_+(k^2),\qquad
 \delta_+(k^2)=\theta(k^0)\delta(k^2).
 \label{eq:cut}
\end{equation}
The energy restriction is part of the physical master definition. It is
retained when evaluating the cut masters, in addition to marking the cuts
in the Kira reduction.

The master measure used in SIDIS is explicitly
\begin{equation}
 M[\mathcal R]=\int\dd^D r\,
 \delta_+(r^2)\delta_+((P_X-r)^2)\,\mathcal R(r).
 \label{eq:cutmeasure}
\end{equation}
There are no absorbed $2\pi$ factors in this definition. The observed $k_1$
is fixed and is not another integrated cut. Thus a tree-level real
contribution becomes a cut one-loop integral for its unresolved phase
space. It does not become a virtual amplitude with an unrestricted loop.

The conversion back to the paper's unobserved-state phase space uses the
normalizations already present in its Eqs.~(34) and (38):
\begin{equation}
 N_{\mathrm{ps}}(D)
 =(2\pi)^D\prod_{a=2,3}(2\pi)^{-(D-1)},\qquad
 N_H=(2\pi)^{-4}.
 \label{eq:normalization}
\end{equation}
The program derives this product and checks it against each channel's
preserved phase-space definition. Dimensional scale factors, charge and
strong-coupling factors, and any remaining spectator weights are inserted
at the documented assembly boundary. They are not part of
Eq.~\eqref{eq:cutmeasure}.

\subsection{Scalar families and Kira reduction}

FeynCalc expresses the contractions as rational loop integrands. Shared
S02 defines candidate families from the actual tree propagator routings,
including $r^2$, $(P_X-r)^2$ and independent uncut denominators.
Shared S03 maps every interference to these families and reconstructs its
original integrand exactly.

Algebraic elimination of dependent denominators is used to construct valid
families. That algebra does not evaluate a phase-space integral. The
integration reduction is performed by Kira, using integration-by-parts
(IBP) relations. Their usual schematic form is
\begin{equation}
 0=\int\dd^D r\,
 \left[\frac{\partial}{\partial r^\mu}
       \bigl(v^\mu\mathcal I(r)\bigr)\right]_{\cut},
 \qquad v\in\{r,p,q,k_1\},
 \label{eq:ibp}
\end{equation}
with the cut-sector rules imposed on the resulting integrals.
Both physical cuts must remain in every surviving master. A pinched cut
sector cannot be treated as an ordinary uncut contribution.

\paragraph{The common scalar coordinates.}
The family construction uses the basis $\{p,q,k_1,r\}$ and solves the
defining scalar-product equations with coordinates
\begin{equation}
 z_1=r^2,\qquad z_2=(P_X-r)^2,\qquad
 a=(p-r)^2,\qquad b=(k_1+r)^2.
 \label{eq:realcoordinates}
\end{equation}
The candidate uncut momentum list is
\begin{equation}
 \{p-r,\ p-(P_X-r),\ k_1+r,\ k_1+(P_X-r),\ q-r,\ q-(P_X-r)\}.
\end{equation}
On the cuts these quadratic denominators are affine in $a,b$. The code
chooses independent pairs by their determinant and checks independence
again in the full $(z_1,z_2,a,b)$ coordinates. The accepted geometry has
twelve candidate families. Each starts with the two physical cuts,
followed by an independent pair of uncut propagators. Every inherited
scalar-product map is checked against its defining equations in the new
coordinates.

\paragraph{Targets and reduction bounds.}
The mapper records integral-index tuples for every numerator and
denominator term. Positive indices specify denominator powers;
nonpositive uncut indices specify absent denominators or numerators.
Kira receives the actual targets, their sectors and preferred masters
retaining both cuts. Reduction bounds are measured from these targets:
the $r$ bound is the largest positive-index sum plus the number of loop
momenta, and the $s$ bound is the largest numerator-index sum. The run
therefore requests the domain needed by the current integrands.

Kira performs initialization, triangular reduction and back substitution,
then exports exact rules. The importer requires every target to have a
rule or be a declared master, unique rule left-hand sides, and complete
termination in the final master inventory. Both physical cuts must have
positive indices in every master. Kira and Fermat operate on symbolic
dimension and invariants throughout this production reduction.

The result is a rational linear combination
\begin{equation}
 R_{c,\Gamma}^{\cut}(\eps;s,t,Q^2,w)
 =\sum_a c_{c,\Gamma a}(D;s,t,Q^2,w)\,M_a(D;s,t,Q^2,w),
 \label{eq:realreduction}
\end{equation}
where $c$ labels a channel and charge sector, and $\Gamma=g,pp$.
The accepted shared reduction has \textbf{315 requested real integrals
and eight cut masters}. These counts are reported by the saved Kira run;
the same master set serves the six channels, with different coefficients.

\subsection{What SubTropica evaluates}

Kira supplies relations between integrals. SubTropica evaluates the
remaining master integrals from explicit parameter representations.
In this workflow it is used as an integration engine, with HyperIntica as
its backend; a master name alone is not treated as a table lookup.
Wolfram Language still coordinates the symbolic work and hosts these
packages.

Shared S08 solves the two cut equations in the $P_X$ rest frame,
selects the positive-energy solution and derives the Jacobian. The saved
tool results are
\begin{equation}
 r^0=|\boldsymbol r|=\frac{\sqrt w}{2},\qquad
 J_{\cut}=2^{-3+2\eps}w^{-\eps},\qquad
 \Omega(n)=\frac{2\pi^{n/2}}{\Gamma(n/2)}.
 \label{eq:cutjacobian}
\end{equation}
Here $J_{\cut}$ is the radial measure after imposing both cuts, and
$\Omega(n)$ is the surface area of the unit sphere in $n$ Euclidean
dimensions. For example, the actual volume-master input is represented as
\begin{equation}
 M_0=J_{\cut}\,\Omega(D-2)
 \int_0^\infty\dd v\,
 [1-c(v)^2]^{-\eps}\left|\frac{\dd c(v)}{\dd v}\right|,
 \qquad c(v)=\frac{1-v}{1+v}.
 \label{eq:volumeinput}
\end{equation}
This is an input to SubTropica, not an inserted evaluated angular table.

On the cuts, every remaining propagator is reconstructed as a linear
scalar product with an effective external vector. The resulting eight
masters fall into four evaluated Euler classes:
\begin{center}
\begin{tabularx}{\textwidth}{@{}l X@{}}
\toprule
Saved class & Parameter integral supplied to SubTropica\\
\midrule
\texttt{cut\_volume} & The positive-energy two-particle cut volume.\\
\texttt{cut\_single\_massive} & One uncut denominator with a non-null effective vector.\\
\texttt{cut\_pair\_one\_massive} & Two denominators, with one null effective vector.\\
\texttt{cut\_pair\_massless} & Two null effective vectors; a separate input before expansion.\\
\bottomrule
\end{tabularx}
\end{center}
The pair inputs use the unexpanded Euler representation referenced in
shared S08~\cite{eulerreference}. Its parameters are obtained from the
current master geometry. The program retains the radial and Gamma-function
prefactors separately and calls \texttt{STIntegrate} on an ordinary
parameter measure. The tuple interface supplies its own internal conversion
to a logarithmic measure; the adapter must not insert that factor twice.

Raw engine output, algebraic-root definitions, physical substitutions and
required $\eps$ orders are saved. The volume class is evaluated through
$\eps^2$, and the other production cut classes through $\eps^1$ in this run.
The required orders are checked against the reduction coefficients and
endpoint needs. An independent normalization check compares two different
representations of the same null-denominator cut integral.

The accepted real route is shared S08/S10/S15 together with the S04 Kira
reduction. The earlier shared S06/S07 uncut-parent route is retained as
history and is not the source of the accepted physical cut-master values.

\subsection{Explicit Euler inputs and their geometric parameters}
\label{sec:euler}

This subsection records the defining parameter integrals constructed in
shared S08. Let $v_i$ be the effective vector reconstructed for an uncut
denominator on the two-cut surface. The program verifies that denominator
against $2r\cdot v_i$ and defines
\begin{equation}
 S_i=P_X\cdot v_i,\qquad
 m_i^2=\frac{w\,v_i^2}{S_i^2},\qquad
 \beta_i=\sqrt{1-m_i^2}.
 \label{eq:effectivevectors}
\end{equation}
Here $m_i^2$ is a dimensionless geometric parameter, not a physical
particle mass. Effective vectors, their powers and all substitutions are
saved for each Kira master. This is how the common classes acquire the
actual $s,t,Q^2,w$ dependence of a channel calculation.

For the single-denominator class the defining density is
\begin{equation}
 M_{\mathrm{single}}=
 \frac{J_{\cut}\Omega(D-2)}{S_1}
 \int_0^\infty\dd v\,
 \frac{[1-c(v)^2]^{-\eps}|c'(v)|}{1-\beta_1 c(v)},
 \qquad c(v)=\frac{1-v}{1+v}.
 \label{eq:singleinput}
\end{equation}
The tuple uses $z_*=(1-\beta_1)/(1+\beta_1)$ as its algebraic parameter.
SubTropica evaluates the parameter integral; the regulated radial and
sphere factors multiply its output afterwards.

For a pair, the source orders the effective vectors so that the second
is null. With denominator powers $j,k$, it defines
\begin{gather}
 y_{12}=\frac{w\,v_1\cdot v_2}{2S_1S_2},\qquad
 z_- =\frac{1-\sqrt{1-m_1^2}}{2y_{12}},\qquad
 z_+ =\frac{1+\sqrt{1-m_1^2}}{2y_{12}},\\
 a_E=j,\qquad b_E=1-k-\eps,\qquad c_E=2-k-2\eps,\\
 I_E=\int_0^1\dd\tau\,
 \tau^{a_E-1}(1-\tau)^{c_E-a_E-1}
 [1-(1-z_-)\tau]^{-b_E}
 [1-(1-z_+)\tau]^{-b_E}.
 \label{eq:paireuler}
\end{gather}
The prefactors stored with this input are
\begin{align}
 \mathcal O_{jk}&=
 2^{2-j-k-2\eps}\pi^{1-\eps}
 \frac{\Gamma(1-k-\eps)}{\Gamma(2-k-2\eps)},\\
 \mathcal E_{jk}&=
 \frac{\Gamma(c_E)}{\Gamma(a_E)\Gamma(c_E-a_E)},\\
 M_{\mathrm{pair}}&=
 \frac{J_{\cut}\mathcal O_{jk}\mathcal E_{jk}}
 {S_1^jS_2^k y_{12}^{\,j}}\,I_E.
 \label{eq:pairprefactor}
\end{align}
These are the unexpanded Euler definitions transcribed from the input
builder. The map $\tau=v/(1+v)$ puts the integral on the positive
half-line used by the tuple. Its Jacobian is included once in the
density. If both vectors are null, the source sets $z_-=0$ in the
defining integrand before expansion and sends a distinct class to the
engine. A limit of an already expanded massive expression is not used
to define that class.

\paragraph{Dimensional orders and algebraic branches.}
An integral's finite term is generally insufficient when its Kira
coefficient has a pole in $\eps$. The input builder measures the Laurent
valuation of the reduction coefficients and keeps the required master
orders. Endpoint assembly also retains the regulated prefactors and
soft powers. The saved root definitions and physical substitutions
accompany every evaluated class so that subsequent logarithms and
dilogarithms are continued on the same branch.

\paragraph{What is reused across channels.}
The master value depends on its cut measure and propagator geometry;
the diagram, color, spin and charge dependence resides in
$c_{c,\Gamma a}$ of Eq.~\eqref{eq:realreduction}. Shared master
evaluation is therefore compatible with different channel coefficients.
It does not license importing a coefficient or average from a different
observed-particle assignment.

\subsection{The accepted cut-master inventory}

The following entries are transcribed from the shared S08 result. Each
tuple lists the four propagator indices in its named family. The first
two entries are the physical cuts. A common class means a common Euler
input; each master still has its own geometric substitutions and
regulated prefactor.

\begin{center}
\begin{tabular}{@{}lll@{}}
\toprule
Family & Indices & Euler class\\
\midrule
R01 & $(1,1,0,0)$ & volume\\
R01 & $(1,1,1,1)$ & pair, both null\\
R02 & $(1,1,1,1)$ & pair, both null\\
R03 & $(1,1,0,1)$ & single, non-null\\
R03 & $(1,1,1,1)$ & pair, one non-null\\
R04 & $(1,1,1,1)$ & pair, one non-null\\
R09 & $(1,1,1,1)$ & pair, one non-null\\
R10 & $(1,1,1,1)$ & pair, one non-null\\
\bottomrule
\end{tabular}
\end{center}

For the powers occurring in this inventory, the saved half-line
densities supplied to \texttt{STIntegrate} are particularly short:
\begin{align}
 f_{\mathrm{volume}}(v)&=
 \frac{2^{1-2\eps}(1+v)^{-2+2\eps}}{v^\eps},\\
 f_{\mathrm{single}}(v;z_*)&=
 \frac{(1+v)^{-1+2\eps}(1+z_*)}{4^\eps v^\eps(v+z_*)},\\
 f_{\mathrm{pair,null}}(v;z_+)&=
 \frac{(1+vz_+)^\eps}{1+v},\\
 f_{\mathrm{pair,non\mbox{-}null}}(v;z_-,z_+)&=
 \frac{[(1+vz_-)(1+vz_+)]^\eps}{1+v}.
 \label{eq:actualcutdensities}
\end{align}
Each density is integrated with $\dd v$ on $[0,\infty)$ and multiplied
by the corresponding saved prefactor. These are the actual parameter
inputs recorded by the tool, not a list of evaluated master values.

\section{Soft limits and distributions}
\label{sec:soft}

A master series at fixed $w>0$ does not yet define its action at $w=0$.
Shared S15 therefore finds the required soft regions before expanding the
regulator. It uses SubTropica's \texttt{STGetRegionVectors} and
\texttt{STPuiseux}, integrates every returned region and keeps the regulated
powers of $w$ until distribution assembly.

The saved plus-distribution convention is
\begin{align}
 L_n(w;B)&=\left[\frac{\log^n(w/B)}{w}\right]_+,
 \qquad 0\leq w\leq B,\label{eq:plus}\\
 \int_0^B\dd w\,L_n(w;B)\,\varphi(w)
 &=\int_0^B\dd w\,\frac{\log^n(w/B)}{w}
       [\varphi(w)-\varphi(0)].
\end{align}
Here $B$ has the same units as $w$. In the paper's convolution variables
its physical upper endpoint is
\begin{equation}
 B=Q^2\left(\frac{1}{\hat x}-1\right)(1-z)
       -\frac{P_{hT}^2}{z}.
\end{equation}
The symbolic coefficient files retain $B$ explicitly.

Shared S17 collects the soft terms of the form
$A_\kappa(\eps)w^{-1-\kappa\eps}$.
The delta coefficient is obtained by evaluating the regulated endpoint
integral in the program; the logarithmic coefficients are obtained from
the exponential series:
\begin{equation}
 E_\kappa(\eps;B)=\int_0^B\dd w\,w^{-1-\kappa\eps},
 \qquad
 \exp[-\kappa\eps\log(w/B)].
 \label{eq:endpoint}
\end{equation}
The calculation uses the convergent orientation $\kappa>0$, $\eps<0$
for this endpoint definition, followed by its Laurent expansion.
It checks cancellation of stronger recoil singularities and of higher
logarithmic plus terms. The saved result separates \texttt{Delta},
\texttt{L0}, \texttt{L1} and the ordinary \texttt{Regular} remainder.

The three channels with virtual corrections retain the physical branches
\begin{equation}
 t=\sigma\omega-s,\qquad \omega>0,\qquad \sigma=+1,-1.
\end{equation}
The key \texttt{0} in the final files is their common boundary $t=-s$,
evaluated by symbolic limits. It must be read from that saved branch;
direct substitution into separately singular terms is not the prescribed
boundary operation.

\subsection{The order of the endpoint operations}

The program first forms the regulated master representation, then finds
its small-$w$ regions. It retains the region scaling exponent together
with each coefficient series. Only after multiplying by the channel's
reduction coefficients and combining the required regions does it
separate delta, plus and ordinary pieces. A fixed-$w$ Taylor expansion
alone would not supply the endpoint action in Eq.~\eqref{eq:plus}.

For each regulated soft row the endpoint contribution is defined by
$A_\kappa(\eps)E_\kappa(\eps;B)$, while the logarithmic coefficients
come from the recorded expansion of
$A_\kappa(\eps)B^{-\kappa\eps}
\exp[-\kappa\eps\log(w/B)]$. The source performs both Laurent
products at the requested order. The ordinary remainder is the expanded
interior expression with its expanded singular soft terms removed.
Consequently its addition to the distribution pieces recovers the same
regulated real integral when applied to a smooth test function.

The endpoint program checks cancellation of powers stronger than the
required $w^{-1-\kappa\eps}$ structure, and its saved coefficient of
$L_2$ vanishes at this order. For Hgg, Hqqbar and Hqqprime the complete
endpoint coefficients vanish in the final assembly. That result is
checked on their own contractions, rather than assumed from the absence
of a Born process.

\subsection{Region enumeration and physical branches}

The accepted source explicitly processes every vector returned by
\texttt{STGetRegionVectors} and calls \texttt{STPuiseux} for each
region. It does not rely on the installed convenience wrapper that, in
the recorded environment, processed only its first returned region.
Raw region outputs and the nonempty regulator series are retained with
the accepted result. This distinction concerns which engine outputs
enter the calculation; it is not an additional analytic approximation.

The branches $s+t>0$ and $s+t<0$ encode the physical continuation of
the combined expression. The boundary entry is obtained from the common
symbolic limit after the relevant terms are combined. A numerical
consumer must use the matching branch and the saved boundary operation,
and must not repair a singular component by assigning it zero.

\section{Virtual integrals and UV renormalization}
\label{sec:virtual}

\subsection{Scalar loop families and fresh master values}

For Hqq, Hqg and Hgq, FeynCalc contracts each one-loop amplitude with its
Born amplitude and expresses the result in complete scalar loop families.
Every mapped integrand is reconstructed before integration. The loop
measure inherited from the generated amplitude is
\begin{equation}
 \frac{\mu^{2\eps}}{(2\pi)^D}\int\dd^D\ell.
\end{equation}
Kira canonicalizes and reduces the virtual targets. The accepted run maps
1625 original targets to 252 canonical targets and seven scalar masters.
The masters use two freshly derived Euler templates, for bubble and box
parents. The evaluated series reach $\eps^1$ for the bubble template and
$\eps^0$ for the box template, as required by their coefficients.

FeynCalc constructs the parameter inputs with
\texttt{FeynmanIntegralPrefactor -> "Unity"}; the generated amplitude's
loop measure remains external. SubTropica evaluates the Euclidean master
series. Shared S19 then applies the recorded physical substitutions and
Feynman prescription to their logarithms and dilogarithms. Only after this
continuation does it form the Hermitian interference,
\begin{equation}
 V_\Gamma=I_\Gamma+I_\Gamma^*,\qquad
 I_\Gamma=\sum_a c^{V}_{\Gamma a}M_a^V.
\end{equation}
Here $I_\Gamma$ includes the stored loop measure and pre-Hermitian
interference convention.

\subsection{Bubble and box parent integrals}

The virtual input builder starts from each master's active propagators,
normalizes the routing signs and constructs the matrix of squared
differences between its loop-momentum shifts. It groups masters by the
zero pattern of that matrix and derives a generic parent Gram matrix.
Every nonzero chord has a symbolic Euclidean parameter $c_{ab}>0$.
The map from that parameter to the actual physical invariant is saved
separately for each master.

FeynCalc produces a homogeneous projective parameter representation.
The program checks its scaling before fixing the final parameter to one.
The two resulting inputs, quoted from shared S12, are
\begin{align}
 I_2(c_{12})={}&i\pi^{2-\eps}\Gamma(\eps)
 \int_0^\infty\dd x_1\,
 \frac{(1+x_1)^{-2+2\eps}}{(c_{12}x_1)^\eps},
 \label{eq:bubbleinput}\\
 I_4(c_{14},c_{23},c_{34})={}&i\pi^{2-\eps}\Gamma(2+\eps)
 \int_0^\infty\dd x_1\dd x_2\dd x_3\,
 \frac{(1+x_1+x_2+x_3)^{2\eps}}
 {(c_{14}x_1+c_{34}x_3+c_{23}x_2x_3)^{2+\eps}}.
 \label{eq:boxinput}
\end{align}
They represent unnormalized Minkowski $\dd^D\ell$ integrals with the
stated FeynCalc prefactor convention. The amplitude's
$\mu^{2\eps}/(2\pi)^D$ factor is applied separately. SubTropica
evaluates these Euler inputs to the orders requested by their measured
coefficients.

The seven physical-master substitutions saved in this run are:
\begin{center}
\small
\begin{tabular}{@{}lll@{}}
\toprule
Kira family and indices & Parent & Physical parameter map\\
\midrule
V01 $(0,1,0,1)$ & $I_2$ & $c_{12}=-s$\\
V02 $(1,1,0,0)$ & $I_2$ & $c_{12}=-t$\\
V03 $(0,1,1,0)$ & $I_2$ & $c_{12}=Q^2$\\
V03 $(1,1,0,0)$ & $I_2$ & $c_{12}=Q^2+s+t$\\
V06 $(1,1,1,1)$ & $I_4$ & $(c_{14},c_{23},c_{34})=(-t,-s,Q^2)$\\
V07 $(1,1,1,1)$ & $I_4$ & $(c_{14},c_{23},c_{34})=(-s,Q^2+s+t,Q^2)$\\
V08 $(1,1,1,1)$ & $I_4$ & $(c_{14},c_{23},c_{34})=(-t,Q^2+s+t,Q^2)$\\
\bottomrule
\end{tabular}
\end{center}
These substitutions accompany the recorded Feynman prescription; their
signs alone are not a replacement for that prescription. In particular,
continuation is performed before adding the conjugate interference.
All 252 canonical targets in eleven families are reduced before the
seven-master assembly. Triangle and tensor structures appearing in
intermediate diagrams are therefore not additional unevaluated outputs.

\subsection{Self energies, scaleless terms and the UV/IR separation}

Self energies and UV residues are also evaluated with Kira and SubTropica.
Shared S14 uses generated quark, gluon/ghost and one-flavor gluon
two-point functions. It isolates the logarithmic large-loop-momentum
terms with an auxiliary mass, reduces them with Kira, and evaluates the
auxiliary vacuum and off-shell bubble masters with SubTropica.
The auxiliary mass is an infrared rearrangement device; it is not a
physical quark mass.

In the stored convention, with the strong couplings removed from the
interference, the accepted common coupling residue is
\begin{equation}
 R_g=-\frac{11N_c-2N_f}{96\pi^2},\qquad
 \Delta_{\MSbar}=\frac1\eps-\gamma_E+\log(4\pi).
 \label{eq:uvresidue}
\end{equation}
This quotes \texttt{CouplingResidue} and \texttt{MSPole} from the new S14
result. The generated Born coupling count gives
\texttt{CouplingWeight}=2 in each of the three Born channels.
The corresponding stored coupling counterterm is
\begin{equation}
 V^{\ct}_{\Gamma,c}
 =n_{g,c}R_g\Delta_{\MSbar}\,B_{\Gamma,c}^{(D)},
 \qquad n_{g,c}=2,
\end{equation}
with the common-regulator and external-field terms included by S19.
The program keeps UV and IR residues separate long enough to check UV
cancellation. Kira's scaleless on-shell boundary and the generated external
field counts fix the external-field bookkeeping. They are not imported
Package-X results. The UV poles cancel before the new renormalized virtual
tensors enter the final real-plus-virtual assembly.

The saved field UV residues in the same stripped-coupling convention are
\begin{align}
 R_q^{\mathrm{UV}}&=-\frac{(N_c-1)(N_c+1)}{32\pi^2N_c},\\
 R_g^{\mathrm{field,UV}}&=-\frac{N_f}{24\pi^2}
                         +\frac{5N_c}{48\pi^2}.
 \label{eq:fieldresidues}
\end{align}
The gluon result records the gauge-sector contribution and the generated
one-flavor fermion contribution separately before multiplying the latter
by $N_f$. These are S14 tool outputs. The coupling residue in
Eq.~\eqref{eq:uvresidue} and the gluon field residue in
Eq.~\eqref{eq:fieldresidues} are different quantities.

For all three Born channels the generated external-field inventory has
two quark fields and one gluon field. S19 constructs the external-field
factor from the product of the square-root field factors and extracts
the interference term at the required order. Its on-shell common-regulator
residue obeys the scaleless boundary returned by Kira. UV and IR
labels are nevertheless kept separate during the cancellation test, so
an on-shell zero does not cause the UV residue needed by renormalization
to be silently discarded.

\subsection{The renormalized virtual assembly}

In the notation of the executed S19 operations, let
$V_\Gamma^{\mathrm{bare}}=I_\Gamma+I_\Gamma^*$ and let
$U_\Gamma$ denote the corresponding Hermitian UV residue. The program
forms the separated expression
\begin{equation}
 V_\Gamma^{\mathrm{sep}}
 =V_\Gamma^{\mathrm{bare}}+
 U_\Gamma\left(\frac{1}{\eps_{\mathrm{UV}}}-\frac{1}{\eps}\right),
 \label{eq:uvseparation}
\end{equation}
then adds the generated external-field interference and the coupling
counterterm, each multiplying the appropriate $D$-dimensional Born
tensor. It requires the coefficient of $1/\eps_{\mathrm{UV}}$ to
vanish. Only after this test are the regulator labels identified and
the remaining series retained for combination with real emission and
factorization. IR poles at this boundary are expected components of
that subsequent combination; a virtual tensor is not required to be
IR finite by itself.

This ordering also fixes the finite $\MSbar$ constants. The pole factor
contains $-\gamma_E+\log(4\pi)$ in the same convention as the loop
measure and the collinear subtraction. A standalone textbook UV
coefficient with a different stripped-coupling or measure convention
cannot be substituted for this complete assembly.

\section{Collinear factorization and finite assembly}
\label{sec:factorization}

\subsection{The two factorization boundaries}

The paper's factorization prescription removes overlaps with the
perturbative parton distributions and fragmentation functions. In its
splitting-kernel convention, Eq.~(46) can be written as
\begin{align}
 \What^{\NLO}_{ij}
 ={}&W^{\NLO}_{ij,\mathrm{unsub}}
 +\frac{\alpha_s}{4\pi}\frac{S_\eps}{\eps}
   \sum_a\int\frac{\dd\eta}{\eta}\,
       \What^{\LO}_{aj}(q,\hat x/\eta,\hat z)P_{a/i}(\eta)
 \nonumber\\
 &+\frac{\alpha_s}{4\pi}\frac{S_\eps}{\eps}
   \sum_b\int\frac{\dd\eta}{\eta^2}\,
       \What^{\LO}_{ib}(q,\hat x,\hat z/\eta)P_{j/b}(\eta).
 \label{eq:factorization}
\end{align}
The integration ranges and endpoint actions are those of the paper's
convolutions. $S_\eps/\eps$ denotes its $\MSbar$ subtraction pole.
The plus signs here follow the negative pole in the perturbative PDF/FF
definitions in the paper. The saved \texttt{Counterterms} or
\texttt{PhaseCT} inputs already include their signs.

The Born tensor inside these convolutions is $D$ dimensional. For an
incoming rescaling, the $p^\mu p^\nu$ contraction also needs the saved
$\eta^{-2}$ projector conversion. The kernel normalizations, fixed-flavor
$P_{qg}$ convention and any quark/antiquark multiplicities are retained
from each channel's own subtraction inputs.

For Hqq, Hqg and Hgq, the loop-independent subtraction results are reused
unchanged. For Hgg, Hqqbar and Hqqprime, shared S16 evaluates the preserved
Born convolution and projector definitions afresh. Neither operation
loads an old integrated real or virtual tensor.

\subsection{Kernels, roots and projector rescaling}

The core subtraction inputs retain the following kernel convention,
quoted here using $v$ for the splitting fraction and
$\mathscr P(v)=[1/(1-v)]_+$:
\begin{align}
 P_{qq}(v)={}&-\frac{(N_c-1)(N_c+1)}{2N_c}
 \bigl[2+2v-3\delta(1-v)-4\mathscr P(v)\bigr],\\
 P_{qg}(v)={}&1-2v+2v^2,\\
 P_{gq}(v)={}&\frac{(N_c-1)(N_c+1)(2-2v+v^2)}{N_c v},\\
 P_{gg}(v)={}&-\frac{1}{3v}\bigl[-12N_c+24N_cv-12N_cv^2+12N_cv^3
 \nonumber\\
 &\hspace{20mm}+2N_fv\delta(1-v)-11N_cv\delta(1-v)
                    -12N_cv\mathscr P(v)\bigr].
 \label{eq:storedkernels}
\end{align}
This notation reproduces the stored kernel definitions without
redistributing their normalization among the coupling prefactor,
species count and convolution. The metadata additionally specifies
kernel-normalization weights. In the Hqq and Hgq core inputs these are
$1/2$ for $P_{qq},P_{gq},P_{gg}$ and $1/4$ for $P_{qg}$, together with
the recorded fermion-species multiplicity 2. Hqg records $1/2$ for its
kernel normalization. These factors belong to those specific contraction
and fixed-flavor inputs. They must be used together with their stored
counterterm definitions, rather than appended to an independently
normalized splitting formula.

The Born delta function fixes the splitting fraction in an interior real
configuration. The accepted core subtraction output records
\begin{gather}
 \eta_{\PDF}=-\frac{t}{w-t},\qquad
 \eta_{\FF}=\frac{s-w}{s},\\
 J_{\PDF}=\frac{1}{w-t},\qquad
 J_{\FF}=\frac{s-w}{s^2},\qquad
 S_{pp,\PDF}=\eta^{-2}.
 \label{eq:ctroots}
\end{gather}
The root and Jacobian are obtained from the defining rescaled Born
kinematics in the source program. The last factor converts the incoming
$pp$ contraction to the original incoming momentum. The metric
contraction has no such projector rescaling. These details matter
independently of the splitting kernel.

Kernel delta and plus terms are transformed with their endpoint action.
The resulting subtraction records carry their own \texttt{Delta},
\texttt{L0} and ordinary pieces, including finite terms from the
$D$-dimensional Born tensor. They are combined with the real and virtual
objects before taking the finite part.

\subsection{Restoring the hard-tensor normalization}

For the three Born channels, the executed assembly has the form
\begin{align}
 \mathbf T_c^{\NLO}
 =\Big[&\mu^{2\eps}N_{\mathrm{ps}}(D)
                \mathbf R_c^{\cut}
       +2\pi\,\mathbf V_c^{\mathrm{ren}}\delta(w)
       +\mathbf C_c^{\mathrm{coll}}\Big]_{\eps^0},
 \label{eq:assembly}\\
 \Fhat_{i,c}^{\NLO}
 &=e_c^2 g_s^4 N_H\,
          \mathcal P_i^{(4)}[\mathbf T_c^{\NLO}],
 \qquad g_s^2=4\pi\alpha_s.
 \label{eq:assemblyprojection}
\end{align}
The bold quantities have $g$ and $pp$ components; $e_c$ here is the fixed
quark charge of the Born process, including Hgq. In Hqq,
$\mathbf R_c^{\cut}$ includes the saved flavor-weighted sum of real sectors.
The stripped-coupling convention in these equations is the one used by
the corresponding S20 core-channel code.

The three channels without Born terms retain their couplings in the
imported real tensors, so S20 uses their own stored normalization instead
of multiplying them by the core-channel prefactor again. Their spectator
weight, phase and collinear term are combined before the same finite
projector operation. The code symbols \texttt{mu} and \texttt{mu2} belong
to the respective saved scale conventions; \texttt{mu2} denotes a squared
scale. Comparisons and numerical consumers use the recorded scale maps.

The finite-part notation in Eq.~\eqref{eq:assembly} is used only after
the program has required
\begin{equation}
 [\eps^{-n}]\,
 \bigl(\mathbf R+\mathbf V+\mathbf C^{\mathrm{coll}}\bigr)=0
 \qquad\text{for every present }n>0,
 \label{eq:poles}
\end{equation}
with all normalization factors included. This check is made separately
for the two contractions and their saved distribution and charge
components. Taking $\eps=0$ before this combination would lose the finite
terms generated by poles multiplying $D$-dependent factors.

\paragraph{The Born contribution.}
For the stripped core Born contractions $\mathbf B_c^{(D)}$, the
executed final projection is
\begin{equation}
 \Fhat_{i,c}^{\LO}
 =e_c^2(4\pi\alpha_s)\,2\pi N_H\,
 \mathcal P_i^{(4)}[\mathbf B_c^{(D)}|_{D=4}]\,\delta(w).
 \label{eq:bornfinal}
\end{equation}
This quotes the production assembly operation. The explicit Born tensors
in the next sections make its channel content unambiguous. The complete
NLO coefficient uses Eq.~\eqref{eq:assemblyprojection}, with its
coupling and charge factors already restored in the final files.

\paragraph{The acceptance boundary.}
No negative Laurent coefficient is discarded merely because a finite
term can be extracted. UV cancellation is required in the virtual route,
then final IR cancellation is required in each contracted, weighted
distribution component. Core boundary branches must agree under their
saved common-limit operation. Only the accepted finite expressions are
exported as the final F hats.

\section{Hqq: an observed quark of the incoming flavor}
\label{sec:hqq}

\subsection{Processes and the Born normalization}

The incoming state is $\gamma^*(q)q_i(p)$ and the tagged parton is
$q_i(k_1)$. The two generated Born diagrams form the QCD Compton
amplitude with recoil $g(k_2)$. Its virtual correction has the same
external state. The core input uses the quark average $A_q$ of
Eq.~\eqref{eq:averages} and removes $e_i^2g_s^2$ from the Born
contractions.

The following equations quote \texttt{BornPg} and \texttt{BornPpp}
from this channel's imported \texttt{s02\_result.wl}. To keep the
polynomial readable, its numerator is displayed separately:
\begin{align}
 \mathcal N_{qq}={}&4Q^4+4Q^2s-2s^2+Ds^2+4Q^2t
                  -8st+2Dst-2t^2+Dt^2,\\
 B_{g,qq}^{(D)}={}&
 \frac{(D-2)(N_c-1)(N_c+1)\mathcal N_{qq}}{2sN_ct},\\
 B_{pp,qq}^{(D)}={}&
 \frac{(D-2)(N_c-1)(N_c+1)(Q^2+s+t)}{2N_c}.
 \label{eq:bornqq}
\end{align}
The stored photon and gluon Ward contractions vanish. These Born tensors
are used both in the LO projection of Eq.~\eqref{eq:bornfinal} and
inside the pole-multiplied NLO terms. Auxiliary Hgq and Hqg Born
contractions are generated for the mixed factorization routes.

\subsection{Real sectors and their flavor weights}

All three real sectors in the process table are required, together with
the virtual correction to $\gamma^*q_i\to q_i g$.
The saved generated tag weights are $1/2$ for the $q_i gg$ sector and
$1$ for each of the same-flavor and distinct-flavor quark sectors.
These weights belong to the imported tensor convention and must not be
applied a second time at final assembly.

The distinct-flavor charge polynomial is retained. Its saved charge
moments have the definition
\begin{equation}
 \mathcal Q_n^{(i)}
 =\sum_{f=1}^{N_f}e_f^n(1-\delta_{if}),\qquad n=1,2.
\end{equation}
The stored observed-flavor label is $i=1$, with the accepted integer domain
$N_f\geq1$. The per-sector flavor weights and these moments accompany the
final payload. The PDF subtraction uses the Hqq Born tensor with $P_{qq}$
and the Hgq Born tensor with $P_{gq}$. The FF subtraction uses the Hqq Born
tensor with $P_{qq}$ and the Hqg Born tensor with the fixed-flavor $P_{qg}$.

More explicitly, the real assembly's charge-component dictionary uses
\begin{equation}
 \mathbf R_{qq}^{\cut}=
 \mathbf R_{qgg,0}+\mathbf R_{\mathrm{same},0}
 +n_{\mathrm{other}}\mathbf R_{\mathrm{distinct},0}
 +\frac{\mathcal Q_1^{(i)}}{e_i}\mathbf R_{\mathrm{distinct},1}
 +\frac{\mathcal Q_2^{(i)}}{e_i^2}\mathbf R_{\mathrm{distinct},2},
 \label{eq:hqqcharges}
\end{equation}
where the saved flavor weight is
\begin{equation}
 n_{\mathrm{other}}=
 \begin{cases}N_f-1,&N_f\geq2,\\0,&\text{otherwise},\end{cases}
 \qquad N_f\in\mathbb Z,\quad N_f\geq1.
\end{equation}
The component subscripts are the stored powers of the other flavor's
charge after removing the fixed Born charge. The overall $e_i^2$ in
Eq.~\eqref{eq:assemblyprojection} restores the physical charge
convention. The moment definitions, weights and domain are part of the
result, so a consumer must not substitute a single value of $N_f$ for
all three charge structures.

Each sector goes through the same two physical cuts and the shared Kira
masters. Its rational coefficients, interference signs and tag weights
remain distinct. The quark-pair sectors include the photon attachments
and their interference before charge summation; neither a flavor count
nor a guessed gluon-splitting term replaces those generated tensors.

\subsection{Subtraction, virtual assembly and final output}

The channel's factorization routes are
\begin{center}
\begin{tabular}{@{}lll@{}}
\toprule
Boundary & Born tensor & Splitting kernel\\
\midrule
Incoming PDF & Hqq & $P_{q/q}$\\
Incoming PDF & Hgq & $P_{g/q}$\\
Outgoing FF & Hqq & $P_{q/q}$\\
Outgoing FF & Hqg & fixed-flavor $P_{q/g}$\\
\bottomrule
\end{tabular}
\end{center}
The incoming $pp$ route uses the projector conversion of
Eq.~\eqref{eq:ctroots}. The virtual contribution uses the channel's
own Born interference and the common new self-energy residues with its
generated field count. All three real sectors, that virtual term and
the four subtraction routes enter the distribution-level pole test.

The final output has an LO delta coefficient and NLO delta, $L_0$,
$L_1$ and ordinary coefficients for both $\Fhat_1$ and $\Fhat_2$.
Its final stage is \path{Hqq/s17_result/}. The two physical branches
and their common boundary retain the charge-moment information of
Eq.~\eqref{eq:hqqcharges}. In hadronic numerics this result is paired
with same-flavor quark or antiquark PDFs and FFs, using the saved
charge-conjugation contract.

\section{Hqg: an observed gluon from an incoming quark}
\label{sec:hqg}

\subsection{Observed-leg assignment and Born tensor}

The observed gluon is $k_1$; the recoil is $q_i(k_2)g(k_3)$ in the real
term and $q_i(k_2)$ in the Born/virtual term. The generated quark average,
charge removal and tagged-real weight are retained. This placement of the
observed leg controls the fragmentation subtraction.
The input kernels are $P_{qq}$ on the PDF side, and $P_{gg}$ and $P_{gq}$
for the gluon and quark parents on the FF side.
The quark charge remains fixed; $N_f$ is the closed massless-flavor count.

The photon and incoming quark are the same species as in Hqq, but the
observed field is now the gluon. FeynArts generates the Born, real and
one-loop amplitudes with this external ordering. The saved Born
contractions, with $e_i^2g_s^2$ removed, are
\begin{align}
 \mathcal N_{qg}={}&-2Q^4+DQ^4+4Q^2s+4s^2-8Q^2t
                  +2DQ^2t+4st-2t^2+Dt^2,\\
 B_{g,qg}^{(D)}={}&-
 \frac{(D-2)(N_c-1)(N_c+1)\mathcal N_{qg}}
 {2sN_c(Q^2+s+t)},\\
 B_{pp,qg}^{(D)}={}&-
 \frac{(D-2)(N_c-1)(N_c+1)t}{2N_c}.
 \label{eq:bornqg}
\end{align}
The source's Born Ward checks pass. These expressions are quoted from
the channel's own accepted Born record. They are not obtained in this
report by exchanging labels in Eq.~\eqref{eq:bornqq}.

\subsection{Real integration and gluon fragmentation}

The generated real process is
\begin{equation}
 \gamma^*(q)q_i(p)\to\obs{g}(k_1)q_i(k_2)g(k_3).
\end{equation}
The cut momenta are $k_2$ and $k_3$. The tagged-square bookkeeping
already accounts for observing a gluon among the generated final
species. The observed gluon is not integrated as a third cut and no
additional spectator symmetry weight is appended to this core tensor.
Kira reduces the contracted real coefficients in the common recoil
geometry, and the physical cut masters supply their integrals.

The factorization table is
\begin{center}
\begin{tabular}{@{}lll@{}}
\toprule
Boundary & Born tensor & Splitting kernel\\
\midrule
Incoming PDF & Hqg & $P_{q/q}$\\
Outgoing FF, gluon parent & Hqg & $P_{g/g}$\\
Outgoing FF, quark parent & Hqq & $P_{g/q}$\\
\bottomrule
\end{tabular}
\end{center}
The last row describes a quark produced in the auxiliary Born process
whose collinear fragmentation produces the observed gluon. This is why
the Hqg subtraction is not the Hqq FF subtraction with an unchanged
kernel label.

\subsection{Virtual term, finite hats and the consumer}

The one-loop correction is interfered with this channel's gluon-tagged
Born amplitude. It uses the seven shared virtual masters and the new
Kira/SubTropica UV input, with two external quark fields and one gluon
field in the recorded count. The renormalized virtual delta term is
combined with the real distribution and the three subtraction routes
before final projection.

The final \path{Hqg/s17_result/} contains both LO and NLO hats with
the core distribution and branch schema. Its fixed quark charge is
supplied by the incoming quark or antiquark in the numerical flavor sum.
The final-state FF is $D_{h/g}$; no final quark-flavor sum is hidden in
that gluon fragmentation factor.

\section{Hgq: an observed fixed-flavor quark from an incoming gluon}
\label{sec:hgq}

\subsection{Photon--gluon fusion and the Born tensor}

The real recoil is $\bar q_i(k_2)g(k_3)$; the Born/virtual recoil is
$\bar q_i(k_2)$. The incoming average uses the generated $D$-dimensional
physical gluon polarization and color convention. It is not replaced by
the incoming-quark average.
The PDF routes use the Hgq Born tensor with $P_{gg}$ and the Hqq Born
tensor with the per-flavor $P_{qg}$. The FF route uses the Hgq Born tensor
with $P_{qq}$. The observed quark charge is fixed, while $N_f$ counts
closed flavors.

The generated Born process is $\gamma^*g\to\obs{q_i}\bar q_i$.
Its saved contractions, with the fixed $e_i^2g_s^2$ factor removed,
are
\begin{align}
 \mathcal N_{gq}={}&-2Q^4+DQ^4-8Q^2s+2DQ^2s-2s^2+Ds^2
                   +4Q^2t+4st+4t^2,\\
 B_{g,gq}^{(D)}={}&\frac{\mathcal N_{gq}}{t(Q^2+s+t)},\\
 B_{pp,gq}^{(D)}={}&\frac{4s}{D-2}.
 \label{eq:borngq}
\end{align}
The Born normalization includes the generated gluon spin and color
average. In particular, the remaining $D$ dependence in
$B_{pp,gq}^{(D)}$ is retained inside the factorization and UV products.
The Born photon and gluon Ward contractions vanish in the accepted input.

\subsection{Real radiation and the subtraction routes}

The real process is
\begin{equation}
 \gamma^*(q)g(p)\to\obs{q_i}(k_1)\bar q_i(k_2)g(k_3).
\end{equation}
The observed quark is fixed before the recoil integration. The incoming
gluon is summed and averaged according to its own physical polarization
projector. The unobserved antiquark and gluon are the two cut lines.
No sum over distinct observed flavors is made inside this fixed-flavor
hard tensor.

The subtraction inputs retain these routes:
\begin{center}
\begin{tabular}{@{}lll@{}}
\toprule
Boundary & Born tensor & Splitting kernel\\
\midrule
Incoming PDF, gluon parent process & Hgq & $P_{g/g}$\\
Incoming PDF, quark Born process & Hqq & per-flavor $P_{q/g}$\\
Outgoing FF & Hgq & $P_{q/q}$\\
\bottomrule
\end{tabular}
\end{center}
The per-flavor normalization and fermion-species metadata of the
incoming mixing term are used together. The auxiliary Hqq tensor is
generated and checked for that contraction convention. Replacing the
whole incoming route by a flavor-summed kernel would change the meaning
of the fixed observed flavor.

\subsection{Virtual correction, output and comparison boundary}

The virtual process has recoil $\bar q_i$ and is the Born--one-loop
interference of photon--gluon fusion. Its loop coefficients differ from
Hqq and Hqg even though the scalar master templates are shared. The new
UV residues enter with the generated external-field count, and its
renormalized term is combined with the cut real result and the three
subtraction routes. The final pole and branch tests precede the
$\Fhat_1,\Fhat_2$ export in \path{Hgq/s17_result/}.

The numerical consumer uses the gluon PDF and sums the fixed-flavor
hard function over active observed quarks and antiquarks with their
charges and FFs. The saved BigTMD coefficient comparison passes within
its stated scope. The inherited MadGraph real-emission sign comparison
is still unresolved; the final IR pole checks do not close that separate
comparison. The numerical figures retain this documented input caveat.

\section{Hgg: an observed gluon from an incoming gluon}
\label{sec:hgg}

\subsection{Generated process and normalization}

This contribution is $\gamma^*g\to gq_f\bar q_f$ at the order calculated.
The unobserved quark-antiquark pair is integrated once. The saved spectator
weight is $1$, and the final charge factor is
\begin{equation}
 \texttt{chargeSum}=\sum_{f=1}^{N_f}e_f^2.
\end{equation}
Quark and antiquark PDF/FF subtraction routes are kept with their own
conventions. Their multiplicities do not introduce a second copy of the
integrated spectator pair. Real integration and factorization provide the
finite result; there is no virtual term at this order.

The generated real amplitude is squared with an incoming-gluon average
and with the gluon $k_1$ treated as observed. The FeynCalc source also
generates the auxiliary quark--gluon Born processes required by the
collinear terms. The real scalar record retains its strong coupling and
charge factors after removing the common electromagnetic coupling used
by the lepton cross section. It explicitly records that the spectator
weight has not yet been applied.

The field list is $\{g(k_1),q_f(k_2),\bar q_f(k_3)\}$. The
unobserved quark and antiquark are different species, giving the saved
$S_{\mathrm{spectator}}=1$. Their two positive-energy constraints form
the cut master. The same shared eight-master basis supplies the integral,
with this channel's gluon-averaged real coefficients.

\subsection{Flavor sum and collinear mixing}

The photon couples to the generated quark line. The fixed representative
charge in that tensor is mapped at final assembly to
$\sum_f e_f^2$. The incoming gluon PDF and outgoing gluon FF do not
perform this electromagnetic sum themselves.

The incoming collinear routes use a quark or antiquark from a gluon
splitting, followed by a Born process with an observed gluon. The
outgoing routes use photon--gluon fusion followed by fragmentation of
the quark or antiquark into the observed gluon. Thus both quark and
antiquark routes are present in the subtraction bookkeeping. Their
species multiplicity is associated with the factorization routes, not
with an extra integration of the same real spectator pair.

\subsection{Finite projection and the real-only normalization}

For this channel and the two following channels, the final code uses the
stored real-only assembly operation
\begin{equation}
 \mathbf T_c^{\mathrm{fin}}=
 \left[S_{\mathrm{spectator},c}(\mu^2)^{2\eps}
 N_HN_{\mathrm{ps}}(D)\,\mathbf R_c^{\cut}
 +\mathbf C_c^{\mathrm{phase}}\right]_{\eps^0}.
 \label{eq:simpleassembly}
\end{equation}
Here $\mu^2$ denotes the code's squared scale \texttt{mu2}, and
$\mathbf C_c^{\mathrm{phase}}$ is the complete stored phase-normalized
Born convolution. The real tensor in this equation retains its own
couplings. The program maps $g_s^2$ to $4\pi\alpha_s$, performs the
channel's charge substitution, and applies $\mathcal P_i^{(4)}$.
The core prefactor of Eq.~\eqref{eq:assemblyprojection} is not applied
again.

The executed Hgg result has zero delta and plus coefficients after
assembly and a finite ordinary NLO coefficient. Both hats are in
\path{Hgg/s11_result/}. In the neutral-pion convolution the associated
luminosity is the gluon PDF times the gluon FF, multiplied by the
charge sum according to the coefficient contract.

\section{Hqqbar: an observed antiquark of the incoming flavor}
\label{sec:hqqbar}

\subsection{Identical spectator quarks}

The process is $\gamma^*q_i\to\bar q_i(k_1)q_i(k_2)q_i(k_3)$.
The incoming flavor and its squared charge, \texttt{eq2}, are fixed.
The two unobserved quarks are identical. S20 reads the generated field
list and computes
\begin{equation}
 S_{\mathrm{spectator}}=\frac{1}{\prod_a n_a!},
\end{equation}
where $n_a$ counts identical unobserved fields. Its saved value for this
channel is $1/2$, applied once because the imported real tensor excludes
that spectator weight.
The initial subtraction uses $q_i\to g$ followed by
$\gamma^*g\to\bar q_i q_i$; the final subtraction uses
$\gamma^*q_i\to gq_i$ followed by $g\to\bar q_i$.
There is no extra flavor sum and no virtual contribution.

Both assignments of the identical quark lines are included in the
generated amplitude before squaring. Their exchange interference and
fermion signs belong to the FeynArts/FeynCalc result. The spectator
factor divides the phase-space counting once; it does not remove an
amplitude or its interference. The source uses a quark incoming average,
while the observed antiquark is summed as the tagged final state.

\subsection{Cut integration and subtraction}

The real tensor is mapped with $k_2=r$, $k_3=P_X-r$ and the same common
invariants as the other channels. The two unobserved quarks are the cut
lines. The rational master coefficients retain the identical-flavor
interference. Kira reduction and SubTropica evaluation then proceed
with the shared physical master inputs.

The incoming PDF route changes the incoming quark into a gluon and uses
the Born photon--gluon-fusion tensor with the antiquark observed. The FF
route uses the Born observed-gluon tensor followed by the fixed-flavor
$g\to\bar q_i$ fragmentation kernel. Shared S16 evaluates these
preserved Born convolution definitions, including their $D$-dimensional
projectors and flavor normalization. There is no virtual amplitude or
UV loop term to add at this order in this channel.

\subsection{Charge convention and final hats}

The spectator factor $1/2$ and the phase normalization enter
Eq.~\eqref{eq:simpleassembly}. The final charge substitution is
\texttt{eq}$^2\to$\texttt{eq2}; \texttt{eq2} denotes the fixed
squared incoming-flavor charge, not a sum over active flavors. The
endpoint tests give zero delta, $L_0$ and $L_1$, and the ordinary finite
coefficients are exported to \path{Hqqbar/s11_result/}.

The hadronic luminosity pairs an incoming quark with the FF of its
antiquark, and includes the conjugate incoming-antiquark process under
the saved consumer convention. Even for a neutral hadron with equal
quark and antiquark FFs, this channel remains a separate hard
coefficient and is counted once in the six-channel sum.

\section{Hqqprime: an observed quark of a different flavor}
\label{sec:hqqprime}

\subsection{The two independent photon charges}

The process is $\gamma^*q_i\to q_j(k_1)q_i(k_2)\bar q_j(k_3)$ with
$j\ne i$. The saved spectator weight is $1$. Both charges are retained
independently, with the polynomial structure
\begin{equation}
 \Fhat_{a,qq'}=e_i^2 A_a+e_i e_j B_a+e_j^2 C_a,\qquad a=1,2.
\end{equation}
This notation labels the saved charge components; no component is dropped
as an integration assumption. The initial subtraction carries $e_j^2$
through the $\gamma^*g\to q_j\bar q_j$ Born process. The final subtraction
carries $e_i^2$ through $\gamma^*q_i\to gq_i$, followed by $g\to q_j$.
The result is for a fixed distinct observed flavor; physical flavor sums
are performed by the consumer.

The generated amplitude retains photon attachment to the incoming-flavor
line and to the distinct-flavor line. Their squares give the two squared
charge structures, and their interference gives the mixed charge
structure. The labels $A_a,B_a,C_a$ refer to those complete coefficient
components after integration and subtraction, rather than to separately
guessed splitting approximations.

\subsection{Real master coefficients and factorization}

The observed $q_j$ is held fixed and the recoil pair $q_i,\bar q_j$
is integrated through the two positive-energy cuts. The distinct flavor
condition makes the spectator species different, as recorded by the
generated field list. The program retains all three charge coefficients
while mapping the real contractions to Kira and assembling their master
values. A mixed-charge coefficient is not dropped because a later
observable may sum it to zero.

There are two factorization routes with different electromagnetic
charges. The PDF route has an incoming splitting $q_i\to g$ followed
by $\gamma^*g\to\obs{q_j}\bar q_j$, so its Born input carries
$e_j^2$. The FF route has $\gamma^*q_i\to\obs{g}q_i$ followed
by $g\to q_j$, so its Born input carries $e_i^2$. The same incoming
and outgoing convolution measures of Eq.~\eqref{eq:factorization}
apply with these fixed-flavor assignments. There is no Born hard part
for the channel itself and hence no NLO virtual interference to add.

\subsection{Final coefficient and flavor summation}

The final assembly uses Eq.~\eqref{eq:simpleassembly} with spectator
weight one. After the charge-resolved pole and endpoint checks, the
output in \path{Hqqprime/s11_result/} is an ordinary finite NLO
coefficient polynomial in the two charges. It is for one fixed ordered
pair of different flavors. Summing over observed species belongs to the
PDF/FF convolution, with an explicit distinct-flavor condition.

For the neutral-pion numerical setup, the consumer verifies equality of
conjugate FFs and the corresponding cancellation of the mixed-charge
luminosity. That is a checked property of this numerical flavor sum.
It neither changes the saved partonic polynomial nor asserts that its
mixed coefficient vanishes for arbitrary fragmentation functions.

\section{Final F hats and the files that contain them}

For Hqq, Hqg and Hgq the output format is
\begin{align}
 \Fhat_{i,c}^{\LO}&=B_{i,c}\,\delta(w),\\
 \Fhat_{i,c}^{\NLO}
 &=D_{i,c}\,\delta(w)+C_{i,c}^{(0)}L_0(w;B)
                  +C_{i,c}^{(1)}L_1(w;B)+R_{i,c}(w).
 \label{eq:finalformat}
\end{align}
The file keys are respectively \texttt{LODelta}, \texttt{Delta},
\texttt{L0}, \texttt{L1} and \texttt{Regular}. The NLO coefficients have
branch keys \texttt{1}, \texttt{-1} and \texttt{0}, as described above.
Their coupling, charge and hard-tensor normalization factors are already
included.

For Hgg, Hqqbar and Hqqprime, the executed endpoint checks give zero
\texttt{Delta}, \texttt{L0} and \texttt{L1}. Their F-hat files contain
the finite NLO ordinary coefficient directly. This is an outcome of the
new calculation and its endpoint checks, not a condition imposed on the
unintegrated real tensors.

\begin{center}
\small
\begin{tabular}{@{}ll@{}}
\toprule
Channel & Full final result relative to this report\\
\midrule
Hqq & \path{Hqq/s17_result/s17_result.wl}\\
Hqg & \path{Hqg/s17_result/s17_result.wl}\\
Hgq & \path{Hgq/s17_result/s17_result.wl}\\
Hgg & \path{Hgg/s11_result/s11_result.wl}\\
Hqqbar & \path{Hqqbar/s11_result/s11_result.wl}\\
Hqqprime & \path{Hqqprime/s11_result/s11_result.wl}\\
\bottomrule
\end{tabular}
\end{center}
Each listed directory also contains the corresponding
\texttt{sNN\_F1\_hat.wl} and \texttt{sNN\_F2\_hat.wl} files.
The full result must accompany those direct expressions because it records
the support, distributions, branches, scale convention and input identities.
The long analytic coefficient expressions are retained in these symbolic
files; the present report explains their construction and interpretation.

\section{Stage map, execution and checks}

Channel-level stages are consecutive links to shared implementations.
The numbers in the following table are \emph{channel} step numbers;
the shared source filenames retain their original execution numbers.

\begin{center}
\small
\begin{tabularx}{\textwidth}{@{}X c c@{}}
\toprule
Operation & Hqq/Hqg/Hgq & Hgg/Hqqbar/Hqqprime\\
\midrule
Import unchanged inputs; define kinematics & s01--s02 & s01--s02\\
Real family mapping and Kira reduction & s03--s04 & s03--s04\\
Physical cut-master inputs; SubTropica evaluation & s06, s08 & s05--s06\\
Real master coefficients; regulated soft regions & s09, s13 & s07--s08\\
Virtual families, reduction, inputs and evaluation & s05, s07, s10--s11 & --\\
New self-energy/UV evaluation & s12 & --\\
Collinear inputs & preserved in s01 & s09\\
Real distribution assembly & s14 & s10\\
Virtual coefficients and renormalized assembly & s15--s16 & --\\
Final pole cancellation and F-hat projection & s17 & s11\\
\bottomrule
\end{tabularx}
\end{center}
The table groups operations by purpose; the numbered channel sequence
and each README give the actual dependencies. Shared integrations and
reductions run once for the applicable channel group.

The accepted calculation ran on Hoffman2 compute nodes. The local folders
retain the main sources, results, logs, Kira databases, SubTropica outputs
and resumable caches. Independent work is divided among a bounded number
of workers with per-operation memory limits, process-tree RSS guards and
persistent job monitors. The recorded accepted runs did not encounter OOM.

The saved gates cover input reconstruction, Born normalization and Ward
identities, complete Kira target coverage, retention of both cuts,
parameter-measure normalization, regulator depth, physical continuation,
UV cancellation, final IR cancellation and the required branch boundaries.
The all-six final acceptance and its execution receipt are in
\path{common/s20_result/}.

The subsequent \path{bigTMD_comparison/} compares the frozen new F hats
with the same reference used by the earlier comparison workflow. Its saved
report gives \textbf{60/60 exact comparisons and 120/120 high-precision
checks passing}. The exact comparison uses the documented rational
reconstruction of printed reference decimals, SU(3), matched charges,
scales and distribution conventions. It is a comparison after production,
and does not set or tune a coefficient. The scope of the public numerical
driver differs from that of the complete reference coefficient files and
is recorded separately.

The inherited \textbf{Hgq MadGraph real-emission sign comparison remains
unresolved}. Neither the passed pole checks nor the BigTMD coefficient
comparison is reported as closing that separate check.

\section{Numerics: from partonic hats to a hadronic cross section}
\label{sec:numerics}

\subsection{Inputs and the quantity being predicted}

The separate \path{numerics/} folder consumes the six new symbolic
results~\cite{localnumerics}. It predicts the neutral-pion transverse-momentum distribution
with the saved H1 bin definitions and cuts. The symbolic hard scattering
ends at the finite F hats; PDFs, FFs, active-flavor choices and
experimental cuts enter here. The accepted numerical run evaluates the
new F hats, rather than loading old channel-integrated bin totals.

The paper's convolution definitions are
\begin{align}
 F_1&=\sum_{ij}\int_x^1\frac{\dd\xi}{\xi}
                    \int_z^1\frac{\dd\zeta}{\zeta^2}
       \Fhat_{1,ij}(x/\xi,z/\zeta,Q^2,q_T)
       f_{i/H}(\xi)D_{h/j}(\zeta),\\
 F_2&=\sum_{ij}\int_x^1\dd\xi
                    \int_z^1\frac{\dd\zeta}{\zeta^2}
       \Fhat_{2,ij}(x/\xi,z/\zeta,Q^2,q_T)
       f_{i/H}(\xi)D_{h/j}(\zeta).
\end{align}
The different $\xi$ measures are part of the paper's $F_1,F_2$
definitions. The fragmentation Jacobian is applied once, and the
distributions act on the complete convolution test function.

The azimuth-independent lepton cross section used by the numerical
definitions is the paper's expression
\begin{equation}
 \frac{\dd\sigma}{\dd x\,\dd y\,\dd z\,\dd p_T^2}
 =\frac{\pi^2\alpha_{\mathrm{EM}}^2}{zxyQ^2}
 \left[xy^2F_1+(1-y)F_2\right].
 \label{eq:leptonxs}
\end{equation}
The maps and weights needed to express it in the integrated variables
are derived in numerical s04 from its defining physics inputs. This
report quotes those saved outputs. It does not introduce a different
lepton normalization or a different $F_2$ convolution measure.

\subsection{Beam energies, cuts and scale choices}

\begin{center}
\begin{tabularx}{\textwidth}{@{}l X@{}}
\toprule
Input & Accepted numerical setting\\
\midrule
Electron and proton energies & $E_e=27.6\,\GeV$, $E_p=820\,\GeV$\\
Inelasticity & $0.1<y<0.6$\\
Photon virtuality panels & $2<Q^2<4.5$, $4.5<Q^2<15$,
 $15<Q^2<70\,\GeV^2$\\
Transverse momentum & H1 bins between $2.5$ and $15\,\GeV$\\
Pion laboratory angle & $5^\circ<\theta_h<25^\circ$\\
Pion laboratory energy & $E_h/E_p>0.01$\\
Common scale & $\mu_R^2=\mu_F^2=\mu_D^2=(Q^2+p_T^2)/2$\\
PDF & MRST2002 NLO, original interpolation, mode 1\\
FF options & KKP neutral pion; Kretzer neutral-pion average\\
Hard-part color & $N_c=3$ in the consumer\\
Mass and azimuth treatment & Massless kinematics; uniform-azimuth
 $F_1,F_2$ approximation\\
\bottomrule
\end{tabularx}
\end{center}

Here $\mu_D$ denotes the fragmentation factorization scale. The
reported plots use the common central-scale choice above. They do not
show an independent scale-variation band. Laboratory cuts are applied
to the pion momentum reconstructed from the saved massless frame map,
not to the partonic recoil or directly to a centre-of-mass polar angle.

\subsection{PDF, FF and coupling conventions}

The native MRST interpolation returns the momentum-weighted PDF
$\xi f_a(\xi,\mu^2)$. The numerical adapter converts it to the
$f_a$ required by the convolution before applying the explicit
$1/\xi$ factor for $F_1$. This separates the PDF library's output
convention from the structure-function measure.

The KKP input uses its saved pion option \texttt{ih=5} and NLO
\texttt{iset=1}. The Kretzer input uses its NLO set and the recorded
neutral combination $(\pi^++\pi^-)/2$. The adapter checks the
neutral quark/antiquark equalities needed by the flavor assembly. It
does not identify quark and antiquark PDFs; those remain distinct.

The retained native flavor order is
\begin{equation}
 \{g,u,\bar u,d,\bar d,s,\bar s,c,\bar c,b,\bar b\}.
\end{equation}
The model supplies the up-type and down-type charges and their
antiparticle signs. The numerical hard parts use four active flavors
below the saved bottom threshold and five above it; the phase-space
contract requires the scale to lie above the charm threshold. The
accepted coupling input uses the EPHOX implicit two-loop convention
with continuity at the bottom threshold and records
\begin{gather}
 \alpha_s(M_Z)=0.1197,\qquad M_Z=91.188\,\GeV,\\
 m_c^2=2.045\,\GeV^2,\qquad m_b^2=18.5\,\GeV^2,\qquad
 \Lambda_4^2=0.1085488995978293\,\GeV^2.
\end{gather}
The electromagnetic input is
$\alpha_{\mathrm{EM}}=0.0072973525643$ and the retained unit
conversion is $1\,\GeV^{-2}=389379372.17185944\,\pb$.
These are the numerical input records, not additional fitted parameters
of the new F-hat calculation.

The PDF and FF adapters are checked on their declared input grids. The
retained PDF domain is $10^{-5}\leq\xi\leq1$ with
$1.25\leq\mu^2\leq10^7\,\GeV^2$, and the FF domain is
$0.01\leq\zeta\leq1$ with
$1\leq\mu^2\leq10^6\,\GeV^2$. Evaluation errors are handled as
errors by the numerical driver; an unsuccessful library call is not
converted to a zero cross-section contribution.

\section{Numerics: recoil variables and laboratory acceptance}
\label{sec:nummap}

\subsection{From the fragmentation fraction to recoil mass}

The integration uses $Q^2,y,p_T,z,\xi,w$ with
\begin{equation}
 x=\frac{Q^2}{4E_eE_p y},\qquad \hat x=\frac{x}{\xi},\qquad
 s=-\frac{Q^2(x-\xi)}{x}.
\end{equation}
The fragmentation fraction is changed to the recoil invariant $w$.
The defining solution recorded in s04, also used in the paper's
convolution map, is
\begin{equation}
 \zeta(w)=
 \frac{\hat x p_T^2+z^2Q^2(1-\hat x)}
 {z[Q^2(1-\hat x)-w\hat x]}.
 \label{eq:zetamap}
\end{equation}
The endpoint value $\zeta_0$ is this map at $w=0$. The FF must be
evaluated separately at $\zeta(w)$ for interior terms and at
$\zeta_0$ for delta terms and endpoint subtraction.

The saved lower bound and recoil endpoint are
\begin{align}
 \xi_{\min}&=
 \frac{x(Q^2z^2-Q^2z-p_T^2)}{Q^2z(z-1)},\\
 B&=Q^2\left(\frac{1}{\hat x}-1\right)(1-z)-\frac{p_T^2}{z},
 \qquad 0\leq w\leq B.
 \label{eq:numbounds}
\end{align}
The source solves the remaining support conditions to give the two
$z$ bounds
\begin{equation}
 z_{\mathrm{low,up}}=
 \frac{\sqrt{Q^2}(x-1)\ \mathbin{\pm}
 \sqrt{(x-1)(Q^2x-Q^2+4p_T^2x)}}
 {2\sqrt{Q^2}(x-1)},
 \label{eq:zlimits}
\end{equation}
where the saved lower root uses the plus sign and the upper root the
minus sign. Configurations with empty physical support do not enter
the convolution. The resulting maps are checked against the invariant
definitions, rather than treating independently sampled $s,t,w$ as
an automatically physical point.

\subsection{Jacobian and invariant evaluation}

The s04 output gives the fragmentation Jacobian explicitly as
\begin{equation}
 J_\zeta(w)=
 -\frac{x\,[Q^2xz^2-Q^2\xi z^2-p_T^2x]}
 {z[Q^2x-Q^2\xi+wx]^2}.
 \label{eq:zetajacobian}
\end{equation}
It is the positive integration Jacobian on the accepted support. All
new coefficient bundles record that this Jacobian is not already
included. The consumer consequently multiplies by $J_\zeta$ once,
along with the separate $1/\zeta^2$ FF measure.

For completeness, the saved invariant maps passed to the hard functions
include
\begin{align}
 t(w)&=\frac{Q^2p_T^2\xi+Q^2wxz^2-p_T^2wx}
 {Q^2xz^2-Q^2\xi z^2-p_T^2x},\\
 u(w)&=-\frac{Q^2\xi z^2(Q^2x-Q^2\xi+wx)}
 {x(Q^2xz^2-Q^2\xi z^2-p_T^2x)},\\
 \hat z(w)&=\frac{z^2(Q^2x-Q^2\xi+wx)}
 {Q^2xz^2-Q^2\xi z^2-p_T^2x}.
\end{align}
These formulas are quoted from the symbolic map result. The branch
selector evaluates $s+t$ at the same point, and the saved boundary
entry is used for its common limit.

The s04 weights multiplying the structure functions after changing
to the integration variables are
\begin{align}
 W_1&=\frac{\pi^2\alpha_{\mathrm{EM}}^2p_T}
 {2E_eE_pQ^2z},\\
 W_2&=-\frac{\pi^2\alpha_{\mathrm{EM}}^2p_T(4E_eE_py-4E_eE_p)}
 {2E_eE_p(Q^2)^2yz}.
 \label{eq:numweights}
\end{align}
The native paired evaluator uses $W_1/\xi$ with $\Fhat_1$ and
$W_2$ with $\Fhat_2$. Their different incoming measures are therefore
preserved explicitly.

\subsection{The pion momentum in the laboratory}

The laboratory acceptance is obtained from the saved massless frame
transformation. With $\phi$ the azimuth used in that map, the output is
\begin{align}
 E_h={}&\frac{2\sqrt{Q^2}p_Tz\sqrt{1-y}\cos\phi+Q^2z^2+p_T^2
                    +yz^2(4E_e^2y-Q^2)}{4E_eyz},\\
 p_{h,x}={}&\sqrt{Q^2}z\sqrt{1-y}+p_T\cos\phi,\\
 p_{h,y}={}&p_T\sin\phi,\\
 p_{h,z}={}&\frac{2\sqrt{Q^2}p_Tz\sqrt{1-y}\cos\phi+Q^2z^2+p_T^2
                    -yz^2(4E_e^2y+Q^2)}{4E_eyz}.
 \label{eq:labmap}
\end{align}
The cuts are imposed on $E_h$ and on the massless polar-angle
definition $\cos\theta_h=p_{h,z}/E_h$. The source solves the energy
and angle inequalities as bounds on $\cos\phi$, clips those bounds
to $[-1,1]$ and calculates the allowed azimuth fraction. An empty
interval gives zero acceptance.

The defining quantity is
\begin{equation}
 \mathcal A_{\mathrm{lab}}(Q^2,y,z,p_T)=\frac{1}{2\pi}
 \int_0^{2\pi}\dd\phi\,
 \mathbf1_{E_h/E_p>0.01}\,
 \mathbf1_{5^\circ<\theta_h<25^\circ}.
 \label{eq:acceptance}
\end{equation}
The accepted numerical source uses the analytic angular fraction
derived from these inequalities, so it does not need a seventh Monte
Carlo variable for $\phi$. It assumes the retained $F_1,F_2$
contribution is uniform in this azimuth. Azimuth-dependent SIDIS
structure functions are outside the plotted approximation; the
acceptance fraction alone does not supply them.

\section{Numerics: assembling all six channel luminosities}
\label{sec:luminosities}

Write $f_a=f_{a/H}(\xi,\mu^2)$ and
$D_a=D_{h/a}(\zeta,\mu^2)$. Let $a$ run over active quarks and
antiquarks, and let $e_a$ carry the charge sign from the model. The
native consumer factors the charge polynomials according to each
export contract and constructs the following luminosities. The formulas
are the definitions used by numerical s08/s09; when the full hat already
contains its charge, the exporter exposes the corresponding normalized
charge component so the factor is not included twice.

\subsection{Hgg, Hgq, Hqg and Hqqbar}

\begin{align}
 \mathcal L_{gg}&=f_gD_g\sum_{f\,\mathrm{quark\ flavors}}e_f^2,\\
 \mathcal L_{gq}&=f_g\sum_a e_a^2D_a,\\
 \mathcal L_{qg}&=D_g\sum_a e_a^2f_a,\\
 \mathcal L_{q\bar q}&=\sum_a e_a^2f_aD_{\bar a}.
 \label{eq:simplelums}
\end{align}
The first sum is over quark flavors, not a doubled quark-plus-antiquark
copy of the Hgg spectator pair. In the other sums quarks and antiquarks
are distinct incoming or observed species. The neutral-FF equalities
are checked where used; they do not erase those species definitions.

\subsection{Hqq charge moments}

Hqq uses same-species luminosities $f_aD_a$ and evaluates its coefficient
with the fixed charge and the other-flavor moments appropriate to that
species. The accepted neutral consumer groups equal-charge species
using its verified charge-conjugation identities. For each group,
\begin{equation}
 \mathcal L_{qq,\mathrm{group}}=\sum_{a\in\mathrm{group}}f_aD_a.
\end{equation}
The hard coefficient is evaluated with its stored
$e_i,\mathcal Q_1^{(i)},\mathcal Q_2^{(i)}$ arguments. The saved
tool-generated values used by the two active-flavor settings are:
\begin{center}
\begin{tabular}{@{}cccc@{}}
\toprule
$N_f$ & Representative charge & $\mathcal Q_1^{(i)}$ & $\mathcal Q_2^{(i)}$\\
\midrule
4 & $2/3$ & $0$ & $2/3$\\
4 & $-1/3$ & $1$ & $1$\\
5 & $2/3$ & $-1/3$ & $7/9$\\
5 & $-1/3$ & $2/3$ & $10/9$\\
\bottomrule
\end{tabular}
\end{center}
These values belong to the numerical flavor model. Production Hqq
continues to retain symbolic moments and its stated $N_f$ domain.

\subsection{Distinct-flavor charge components}

For Hqqprime, the pair $(a,b)$ is ordered and its quark-flavor names
must differ, independently of whether a species is a quark or an
antiquark. The source constructs
\begin{align}
 \mathcal L_{\mathrm{incoming}}&=
 \sum_{\mathrm{flavor}(a)\ne\mathrm{flavor}(b)}e_a^2 f_aD_b,\\
 \mathcal L_{\mathrm{prime}}&=
 \sum_{\mathrm{flavor}(a)\ne\mathrm{flavor}(b)}e_b^2 f_aD_b,\\
 \mathcal L_{\mathrm{mixed}}&=
 \sum_{\mathrm{flavor}(a)\ne\mathrm{flavor}(b)}e_ae_b f_aD_b.
 \label{eq:primelums}
\end{align}
They multiply the corresponding three charge coefficients of
Section~\ref{sec:hqqprime}. The neutral-pion input passes the symbolic
and native mixed-luminosity cancellation checks. The mixed partonic
coefficient remains in the symbolic result. No factor based only on
the number of active flavors replaces these ordered sums.

\section{Numerics: distributions, precision and bin integration}
\label{sec:binintegration}

\subsection{The complete endpoint test function}

For a distribution component $d$ and a channel/charge component $c$,
the paired hard evaluator forms the saved combination
\begin{equation}
 K_{c,d}(w)=\frac{J_\zeta(w)}{\zeta(w)^2}
 \left[\frac{W_1}{\xi}\Fhat_{1,c,d}(w)
                         +W_2\Fhat_{2,c,d}(w)\right].
 \label{eq:pairedkernel}
\end{equation}
The channel's compatible charge luminosity is
$\mathcal L_c(w)$, using the FF at $\zeta(w)$. In this notation
$G_{c,d}(w)=K_{c,d}(w)\mathcal L_c(w)$ names the complete test-function
product used by the native source. The product includes the kinematic
dependence of the coefficient, its Jacobian, the FF and the PDF, rather
than subtracting only the value of the analytic F-hat expression.

For a core channel the implemented recoil integration is represented by
\begin{align}
 \mathcal I_{c}^{\LO}={}&G_{c,\mathrm{Born}}(0),\\
 \mathcal I_{c}^{\NLO}={}&G_{c,\Delta}(0)
 +\int_0^B\dd w\,\left\{
 G_{c,\mathrm{Regular}}(w)
 +\frac{G_{c,L_0}(w)-G_{c,L_0}(0)}{w}\right.
 \nonumber\\[-1mm]
 &\hspace{39mm}\left.
 +\frac{G_{c,L_1}(w)-G_{c,L_1}(0)}{w}\log\frac{w}{B}
 \right\}.
 \label{eq:numericalplus}
\end{align}
This is the numerical source's implementation of the saved distribution
convention. Its delta terms are evaluated at the endpoint without a
recoil integration Jacobian. The ordinary and endpoint-subtracted
interior terms carry that Jacobian. Hgg, Hqqbar and Hqqprime enter with
their ordinary NLO terms and no LO contribution.

\subsection{Six-dimensional sampling map}

Each accepted bin uses six unit-cube variables. The transformations are
logarithmic for $Q^2,p_T,z,\xi$ and linear for $y,w$, with the
kinematic lower and upper limits evaluated in their dependency order.
The definitions in the source are
\begin{align}
 x_{\mathrm{lin}}(v)&=a+v(b-a),&
 J_{\mathrm{lin}}&=b-a,\\
 x_{\log}(v)&=a\exp[v\log(b/a)],&
 J_{\log}&=x_{\log}(v)\log(b/a),\qquad 0\leq v\leq1.
\end{align}
At a sampled point the outer weight contains the five non-recoil
Jacobians, $\mathcal A_{\mathrm{lab}}$ and the pb conversion. The
linear $w$ Jacobian multiplies only the interior terms in
Eq.~\eqref{eq:numericalplus}. The endpoint contribution is independent
of the sixth coordinate and is included with its unit-cube measure.

The PDF is evaluated once for the sampled $\xi,\mu^2$ and used in
all channel components. FFs are evaluated at both the interior and
endpoint fragmentation fractions. All six channels and both FF choices
use the same sample points, retaining the covariance of their sum.
Empty physical support or failed experimental cuts produce a defined
zero contribution. A failed hard-function evaluation is an error and
is not treated as an empty phase-space point.

\subsection{Stable evaluation of symbolic coefficients}

The symbolic export retains all distribution, branch, coupling, scale
and charge contracts. The numerical build contains 110 scalar
coefficients arranged into 55 $F_1,F_2$ pairs. The native evaluator
uses the recorded multiprecision complex arithmetic, with the
100-decimal-digit Boost implementation and a high-precision
dilogarithm path. Opposite large terms are combined at the prescribed
precision before returning the real observable contribution.

Two saved comparison stages validate that translation. Numerical s07
records 7662 native/Wolfram checks. Numerical s10 records 1274
independent checks using 100- and 140-digit evaluation. These are
coefficient-evaluation tests against the new symbolic inputs; they are
separate from the statistical convergence of an integrated bin.

\subsection{VEGAS iterations and statistical acceptance}

The run uses VEGAS 6.3 with a recorded seed for every bin. Its initial
configuration has three adaptation iterations, then eight production
iterations with 1000 requested evaluations per iteration and adaptation
disabled. If the acceptance conditions fail, the driver increases the
requested evaluations by the recorded doubling rule, up to its bounded
refinement limit. Accepted outputs preserve the actual evaluation
counts and iteration records.

The final estimator uses the arithmetic average of production
iterations, as recorded by \texttt{weighted=False}. The joint result
retains all 24 signed components: two FF options, six channels and LO/NLO
orders. Their covariance is used for the uncertainty of the total;
signed channel contributions are not converted into positive weights.
The weighted iteration estimate is used only for its consistency
diagnostic, not silently substituted for the reported estimator.

The recorded error requirement for each FF prediction is
\begin{equation}
 \delta\sigma\leq
 \max\!\left(0.05|\sigma|,
 0.005\,\left.\frac{\dd\sigma}{\dd p_T}\right|_{\mathrm{data}}
 \Delta p_T\right).
 \label{eq:mcgate}
\end{equation}
Additional gates require finite covariance, consistent nonnegative
variances, agreement of the first and second iteration halves within
the saved three-error pull, a chi-square tail probability at least
$0.001$, and reconstruction of each FF total from its channel
components. All 17 accepted bins pass their saved gates. Those gates
establish numerical integration quality for the specified setup; they
are not a test of agreement with the experimental central value.

The plotted quantity is the bin integral divided by its transverse
momentum width,
\begin{equation}
 \left\langle\frac{\dd\sigma}{\dd p_T}\right\rangle_b
 =\frac{\sigma_b}{\Delta p_{T,b}}.
 \label{eq:binaverage}
\end{equation}
The integration error is divided by the same width. The calculation
integrates over the full $Q^2,y,p_T$ bin and over the convolution
variables with acceptance, rather than evaluating the hard function
at one representative point.

\section{Numerical results, plots and comparison conventions}
\label{sec:numresults}

\subsection{The saved bin averages}

Table~\ref{tab:bins} transcribes the saved SIDIS plot result. H1
uncertainties are the recorded total experimental errors. KKP and
Kretzer uncertainties are numerical integration errors only. Values
are rounded for display; the JSON results retain the accepted precision.
All cross-section entries in the table have units $\pb/\GeV$.

% BEGIN GENERATED BIN TABLE
\begingroup
\small
\setlength{\tabcolsep}{4pt}
\begin{longtable}{@{}llrrr@{}}
\caption{Accepted SIDIS bin averages and saved H1 data.\label{tab:bins}}\\
\toprule
$Q^2\ [\GeV^2]$ & $p_T\ [\GeV]$ & H1 & SIDIS KKP & SIDIS Kretzer\\
\midrule
\endfirsthead
\toprule
$Q^2\ [\GeV^2]$ & $p_T\ [\GeV]$ & H1 & SIDIS KKP & SIDIS Kretzer\\
\midrule
\endhead
\bottomrule
\endfoot
2--4.5 & 2.5--2.8 & $215\pm 36$ & $199.9\pm 3.47$ & $142.6\pm 2.71$\\
 & 2.8--3.3 & $146\pm 27$ & $130.1\pm 3.19$ & $95.75\pm 3.08$\\
 & 3.3--4 & $60\pm 14$ & $64.9\pm 1.32$ & $47.29\pm 1.03$\\
 & 4--5.2 & $21.7\pm 4.7$ & $24.31\pm 0.34$ & $16.94\pm 0.254$\\
 & 5.2--8 & $4.37\pm 0.83$ & $5.264\pm 0.0767$ & $3.665\pm 0.0526$\\
 & 8--15 & $0.29\pm 0.26$ & $0.4058\pm 0.00633$ & $0.2894\pm 0.00465$\\
\midrule
4.5--15 & 2.5--2.8 & $200\pm 34$ & $187.7\pm 3.18$ & $139.3\pm 2.55$\\
 & 2.8--3.4 & $122\pm 21$ & $120.9\pm 1.81$ & $88.54\pm 1.42$\\
 & 3.4--4.1 & $58\pm 11$ & $64.04\pm 1.06$ & $45.91\pm 0.801$\\
 & 4.1--5.2 & $23\pm 4.9$ & $26.31\pm 0.441$ & $18.7\pm 0.338$\\
 & 5.2--8 & $5.8\pm 1.2$ & $6.189\pm 0.108$ & $4.409\pm 0.0859$\\
 & 8--15 & $0.31\pm 0.12$ & $0.5181\pm 0.00849$ & $0.3713\pm 0.00608$\\
\midrule
15--70 & 2.5--2.9 & $118\pm 18$ & $100.3\pm 2.6$ & $77.83\pm 2.41$\\
 & 2.9--3.5 & $69.1\pm 8.8$ & $67.78\pm 1.32$ & $52.85\pm 1.23$\\
 & 3.5--4.7 & $27.7\pm 4.7$ & $30.69\pm 0.424$ & $23.65\pm 0.339$\\
 & 4.7--8 & $5.41\pm 0.96$ & $6.817\pm 0.122$ & $5.125\pm 0.0962$\\
 & 8--15 & $0.21\pm 0.06$ & $0.4823\pm 0.00618$ & $0.3536\pm 0.00454$\\
\end{longtable}
\endgroup
% END GENERATED BIN TABLE

The corresponding accepted integrated cross sections are given in
Table~\ref{tab:integrals}. The final row sums the measured bin domains
included in these three panels; it is not an extrapolation outside
their cuts and transverse-momentum range.

% BEGIN GENERATED INTEGRAL TABLE
\begin{table}[htbp]
\centering
\caption{Integrated cross sections in pb, with integration errors.\label{tab:integrals}}
\begin{tabular}{@{}lrr@{}}
\toprule
$Q^2\ [\GeV^2]$ & SIDIS KKP & SIDIS Kretzer\\
\midrule
2--4.5 & $217.2\pm 2.17$ & $156.4\pm 1.92$\\
4.5--15 & $223.6\pm 1.72$ & $162.6\pm 1.35$\\
15--70 & $143.5\pm 1.46$ & $110.6\pm 1.32$\\
\midrule
All included bins & $584.3\pm 3.13$ & $429.6\pm 2.69$\\
\bottomrule
\end{tabular}
\end{table}
% END GENERATED INTEGRAL TABLE

\subsection{What the plotted comparisons mean}

The solid blue and orange predictions in Figure~\ref{fig:crosssection}
use the new SIDIS hats with KKP and Kretzer fragmentation, respectively.
They are LO plus NLO bin averages with all six channels. The black
points are the saved H1 measurements with total experimental errors.
The complete BigTMD predictions use the recorded full coefficient
formulas and the matched numerical inputs. Only the full six-channel
BigTMD variant is present in these SIDIS figures.

The dashed paper curves are the preserved original EPS curve vertices
identified by the plotting inputs: hep-ph/0411212, Fig.~4, and
arXiv:1808.04396, Fig.~3. They are displayed as curves. Their values
used in the ratio figure are interpolated at each bin centre, whereas
the SIDIS and BigTMD predictions are bin averages. A curve-to-bin-centre
ratio is therefore not a reconstruction of the paper's bin-integrated
prediction, particularly for the broadest transverse-momentum bins.

The hard-coefficient BigTMD comparison in
\path{bigTMD_comparison/} and the plotted cross-section overlay answer
different questions. The former compares partonic coefficient formulas
under matched symbolic conventions. The latter includes PDFs, FFs,
phase-space integration, cuts and numerical errors. The saved exact
coefficient agreement is not a claim that independent Monte Carlo
central values must coincide digit for digit.

\subsection{Ratio definition and uncertainties}

The source defines the plotted theory-to-data ratio and its integration
error by
\begin{equation}
 r_b=\frac{\langle\dd\sigma/\dd p_T\rangle_b^{\mathrm{theory}}}
 { (\dd\sigma/\dd p_T)_b^{\mathrm{H1}}},\qquad
 \delta r_b^{\mathrm{int}}=
 \frac{\delta\langle\dd\sigma/\dd p_T\rangle_b^{\mathrm{int}}}
 { (\dd\sigma/\dd p_T)_b^{\mathrm{H1}}}.
 \label{eq:ratios}
\end{equation}
The gray band in Figure~\ref{fig:ratios} is the H1 relative total
uncertainty around one. It is shown separately from the theory
integration errors. The plots do not include PDF uncertainty, FF
uncertainty, scale variation or an uncertainty from the azimuth
approximation. The inherited Hgq MadGraph comparison remains a stated
input caveat, as described in Section~\ref{sec:hgq}.

\clearpage
\begin{figure}[p]
 \centering
 \includegraphics[width=\textwidth,height=0.83\textheight,keepaspectratio]
 {numerics/dsigmapibydpt/s03_cross_section.pdf}
 \caption{The saved neutral-pion LO+NLO transverse-momentum comparison
 using the six new SIDIS channels. Panels show the three $Q^2$ ranges.
 SIDIS and full BigTMD predictions are bin averages; paper inputs are
 curves. Theory errors represent numerical integration only. The
 setup and inherited Hgq caveat are stated in the text.}
 \label{fig:crosssection}
\end{figure}
\clearpage
\begin{figure}[p]
 \centering
 \includegraphics[width=\textwidth,height=0.83\textheight,keepaspectratio]
 {numerics/dsigmapibydpt/s03_ratios.pdf}
 \caption{The saved theory-to-H1 ratios with KKP and Kretzer FFs.
 SIDIS and full BigTMD entries use bin averages; paper ratios use
 interpolated bin-centre values. The gray band shows experimental
 relative total errors, while theory bars show integration errors.
 Ratios use the definition in Eq.~\eqref{eq:ratios}.}
 \label{fig:ratios}
\end{figure}
\clearpage

\section{Files, reproducibility and the boundary of this report}

\subsection{Symbolic channel and shared artifacts}

Each channel directory contains its sequential stage sources and the
corresponding \texttt{sNN\_result} folders. Its README records the
channel conventions, normalization, accepted artifact dependencies and
consumer boundaries. The \path{common/} directory holds implementations
and results shared by channels, including Kira reduction databases,
SubTropica inputs, evaluated master series and the final acceptance
receipt. The layout adapter maps the original execution paths to the
organized local folders; its preserved source and relocation records
identify that path conversion.

The \path{.cluster/} directory contains the cluster execution support.
The accepted production uses bounded workers and resumable caches,
keeping the main source and result artifacts locally. Sharing one
master reduction across the channels saves repeated integration work
without merging their charge or spectator conventions. The current
document expansion does not rerun those physics stages.

\subsection{Numerical artifact sequence}

The numerical README specifies the stage contracts. The sequence starts
with frozen hard-function inputs and exports their branch and charge
components. It then derives the physical kinematic maps, builds the
native multiprecision evaluator, validates its PDF/FF interfaces and
coefficient values, assembles flavor luminosities and integrates bins.
The independent precision comparison precedes publication of the
accepted plot bundle. The following records connect the equations and
plots in this report to the retained outputs:

\begin{center}
\small
\begin{tabularx}{\textwidth}{@{}l X@{}}
\toprule
Path within \texttt{numerics/} & Content used here\\
\midrule
\path{s04_result} & Defining physics inputs and symbolic maps,
 Jacobians, lepton weights and laboratory momentum.\\
\path{s06_result} & Numerical PDF, FF and coupling conventions
 and interface checks.\\
\path{s08_result} & Flavor, charge, neutral-FF and acceptance
 contracts used in the convolution.\\
\path{s09_result} & Seventeen accepted bins, components,
 covariance, iteration records and integrated cross sections.\\
\path{s10_result} & Independent coefficient precision checks.\\
\path{s11_result} & Accepted numerical aggregate and its output
 identity.\\
\path{dsigmapibydpt/s03_result} & Bin averages, ratios, comparison
 inputs and figure identities.\\
\bottomrule
\end{tabularx}
\end{center}

The accepted numerical run records no OOM and a peak memory of about
$1.48$ GiB. Reusing its finished figures in this report requires only
PDF typesetting. The source and regenerated report PDF are stored
together in the SIDIS directory.

\subsection{What an extension would have to supply}

This is an NLO result. Reverse unitarity, Kira and explicit master
evaluation provide a reusable integration organization, but this run
does not establish an NNLO result. An extension would need its own
higher-order amplitudes, cut and virtual families, master integrals and
regulator depths, renormalization and factorization inputs, and final
cancellation checks. None of those missing higher-order objects is
supplied merely by the availability of the current two parent templates
or by the agreement of the NLO coefficients with a reference.

\begin{thebibliography}{9}
\raggedright

\bibitem{sidispaper}
B. Wang, J. O. Gonzalez-Hernandez, T. C. Rogers and N. Sato,
\emph{Large Transverse Momentum in Semi-Inclusive Deeply Inelastic
Scattering Beyond Lowest Order}, arXiv:1903.01529v1 (2019).
Supplied local paper in \path{scripts/}; especially Eqs.~(9), (12),
(16)--(19), (21)--(24), (32), (34), (38) and (46)--(53).

\bibitem{standalone}
Supplied \path{standalone_nlo_qcd_reference/standalone_nlo_qcd-main/},
including \path{manual/standalone.tex} and its FeynCalc workflow.
This supplies the reverse-unitarity/Kira/SubTropica integration approach.
The actual SIDIS routing, physical cut representations and normalization
are defined by the SIDIS programs.

\bibitem{eulerreference}
Somogyi, arXiv:1101.3557, Eqs.~(7), (56) and Appendix B.
This is the unevaluated Euler-representation reference recorded in
\path{common/s08_cut_master_inputs.wl} and its accepted result.
SubTropica evaluates the SIDIS inputs built from that representation.

\bibitem{localcalculation}
Local implementation and accepted results:
\path{common/s02_definitions.wl},
\path{common/s04_reduce_real.wl},
\path{common/s08_cut_master_inputs.wl},
\path{common/s10_evaluate_cut_masters.wl},
\path{common/s14_uv_residues.wl},
\path{common/s17_assemble_real.wl},
\path{common/s19_assemble_virtual.wl} and
\path{common/s20_final_hats.wl};
the six channel \path{README.md} files; and the corresponding numbered
result directories. Original executed sources and the layout manifest
are retained under \path{common/previous_runs/layout_sources/} and
\path{common/s21_result/layout.json}.

\bibitem{localnumerics}
Local numerical implementation and accepted outputs:
\path{numerics/README.md};
\path{numerics/s04_derive_convolution_maps.py},
\path{numerics/s06_prepare_pdf_ff.py},
\path{numerics/s08_derive_consumer_actions.py},
\path{numerics/s09_convolve_channels.py},
\path{numerics/s10_validate_precision.py},
\path{numerics/s11_aggregate_replacement.py};
their numbered result files; and
\path{numerics/dsigmapibydpt/s03_result},
which records the H1 points, paper-curve inputs, full BigTMD overlay,
ratios and figure identities used in this report.

\end{thebibliography}
\end{document}
